
\documentclass[UTF8,11pt,a4paper]{ctexart}
\usepackage[a4paper,left=23mm,right=23mm,top=24mm,bottom=24mm]{geometry}
\usepackage{fontspec}
\setmainfont{DejaVu Sans}
\setCJKmainfont{Noto Serif CJK SC}
\usepackage{amsmath,amssymb}
\usepackage{xcolor}
\usepackage[colorlinks=true,linkcolor=blue!45!black,unicode=true]{hyperref}
\usepackage{enumitem}
\usepackage{needspace}
\usepackage{fancyhdr}
\setlength{\parindent}{2em}
\setlength{\parskip}{0.25em}
\setlist[enumerate]{leftmargin=2.6em,itemsep=0.4em,topsep=0.5em}
\pagestyle{fancy}
\fancyhf{}
\fancyhead[L]{\small 三千弱水 · 大学化学高阶综合题库}
\fancyhead[R]{\small 待专家审校稿}
\fancyfoot[C]{\thepage}
\setlength{\headheight}{15pt}
\sloppy

\begin{document}
\begin{center}
{\LARGE\bfseries 大学化学高阶综合题 · 第5辑 (50题)}\\[0.8em]
{\large 题目、分步解析与数值核验}\\[0.5em]
{\small \textcolor{red!65!black}{审校稿：未经化学专家核准；禁止直接视为正式试卷或出版物}}
\end{center}
\vspace{1em}
\noindent 本册从仓库结构化 JSONL 自动生成。增补和修改应只编辑源题目数据，并重新执行统一编译。
\par 机械方程残差检查只能说明指定模型的内部数值一致，不保证模型现实可行性。第六辑另须完成跨辑模型去重审查。
\tableofcontents
\clearpage
\section*{第5辑 · 题目及解答}
\addcontentsline{toc}{section}{第5辑}
\subsection*{无机化学}
\addcontentsline{toc}{subsection}{第5辑 / 无机化学}
\Needspace{7\baselineskip}
\subsubsection*{UC-ADV5-I20 · 热激发Frenkel空位—间隙对的自洽浓度与形成焓}
\noindent\textbf{题目}\par
在具有等数目可占据正常阳离子位和间隙位的假想AB晶体中，Frenkel缺陷对的平衡为A\_A⇌V\_A+A\_i；定义同一标准位点分数x为每种缺陷的分数，空位生成对的理想平衡常数K=x²/(1−x)。900 K时测得x=0.002000，缺陷对形成焓ΔH=+40.00 kJ/mol且视为常数。求900 K的K、1100 K的K及x、缺陷密度相对增量，并证明自洽正根唯一。不能在计算中忽略正常位占据率1−x却又声称精确。\par
\medskip\noindent\textbf{分步解答}
\begin{enumerate}
\item 每个Frenkel对产生一个正常位空位与一个间隙离子，因此[V]=[A\_i]=x，正常位占据1−x。
\item 900K的质量作用常数为K1=x1²/(1−x1)=0.000004008016032。
\item van’t Hoff关系ln(K2/K1)=−ΔH/R(1/T2−1/T1)，代入得1100K常数K2=0.00001059300823。
\item 由x2²=K2(1−x2)，得x²+K2 x−K2=0，唯一正根x2=[−K2+√(K2²+4K2)]/2=0.003249397995。
\item 缺陷分数相对增长x2/x1=1.624698997；因为ΔH>0，升温使正向缺陷生成有利。
\item 函数x²/(1−x)对0<x<1严格递增，从0增至∞，所以求得的正根唯一，而另一二次根为负。
\item 数值逆核验：按独立化学计量、平衡和约束关系复算，三个无量纲检验残差分别为 3.198e-16、0.000e+0、0.000e+0；满足题定理想模型，但不代表专家审定。
\end{enumerate}
\noindent\textbf{答案摘要}\par
K900 = 4.00801603e-06 dimensionless；K1100 = 1.05930082e-05 dimensionless；x1100 = 0.00324939799 site fraction；relative = 1.624699 ratio\par
\noindent\textbf{模型前提：}
理想稀缺Frenkel位点对，缺陷无聚集且两个子晶格标准态一致，形成焓常数。\par
\medskip\hrule\medskip
\Needspace{7\baselineskip}
\subsubsection*{UC-ADV5-I21 · 三价金属二聚水解的单体—二聚体分布反演}
\noindent\textbf{题目}\par
一份100.0mL酸性水样中假想三价金属以单核M3+及二核M2(OH)2\textasciicircum{}4+形式存在。条件二聚反应为2M3++2H2O ⇌ M2(OH)2\textasciicircum{}4++2H+，pH由强缓冲固定，已将H+效应吸收到K\_cond=[dimer]/[M3+]²=450 L/mol。金属元素总分析浓度CM=0.01200M（以单个M原子计），无沉淀和其它物种。求单体、二聚体浓度、二聚金属原子比例、单体与二聚物摩尔数及二聚程度随K增大方向。\par
\medskip\noindent\textbf{分步解答}
\begin{enumerate}
\item 设二聚体的分子浓度x；总金属元素浓度C=[M3+]+2[dimer]，所以[M3+]=C−2x。
\item 质量作用关系x=K(C−2x)²，物理区间0<x<C/2。
\item 在该区间x/(C−2x)²严格递增，所以物理解唯一，可用二分法求x=0.00443102554 M。
\item 游离单体浓度m=0.00313794892 M，二聚体浓度d=0.00443102554 M，注意一个二聚体贡献两个M原子。
\item 二聚金属原子分数2d/C=0.7385042567，100mL水样单体mol=0.000313794892、二聚体分子mol=0.000443102554。
\item 增加K使方程左侧的目标值增大，由单调性可知x增大；不能把dimer浓度直接当作总金属分析浓度。
\item 数值逆核验：按独立化学计量、平衡和约束关系复算，三个无量纲检验残差分别为 2.220e-16、0.000e+0、0.000e+0；满足题定理想模型，但不代表专家审定。
\end{enumerate}
\noindent\textbf{答案摘要}\par
monomer = 0.00313794892 mol/L；dimer = 0.00443102554 mol/L；atom\_fraction = 0.738504257 fraction；dimer\_mol = 0.000443102554 mol\par
\noindent\textbf{模型前提：}
条件形成常数已包含固定pH与水的活度，二核形式只作为假想二聚模型，忽略真实多核水解的竞争。\par
\medskip\hrule\medskip
\Needspace{7\baselineskip}
\subsubsection*{UC-ADV5-I22 · 配体质子化与单核配位的竞争质子账本}
\noindent\textbf{题目}\par
在pH=9.000的外部缓冲水溶液中，一元单负配体L−的共轭酸HL满足pKa=8.000。假想三价金属M3+总浓度0.003000M，仅形成ML2+，条件形成常数对游离L−给β=10\textasciicircum{}3.2 L/mol。游离酸碱池[HL]+[L−]与络合消耗的L之和为L总=0.01200M。求自由L−、HL、游离M3+、ML2+和束缚配体占比；给出唯一根证明并核验全部配体和金属守恒。\par
\medskip\noindent\textbf{分步解答}
\begin{enumerate}
\item 固定pH−pKa=1，故游离酸碱形态的比[HL]/[L−]=10\textasciicircum{}(pKa−pH)=0.1000。设l=[L−]，则[HL]=0.1l。
\item 配位平衡给[ML2+]=βl[M3+]；金属元素守恒CM=[M3+]+[ML2+]，由此[M3+]=CM/(1+βl)。
\item 配体物料守恒CL=l+0.1l+[ML2+]=1.1l+CMβl/(1+βl)；酸碱池与配位池都要计入。
\item 左侧在l≥0严格单调递增，且l=0时为0，l=CL时大于CL，故在(0,CL)有且仅有一个正根。
\item 代入β=10\textasciicircum{}3.2求l=0.008372934728 mol/L，HL=0.0008372934728 mol/L。
\item 平衡游离M3+=0.0002102282011 mol/L，络合ML2+=0.002789771799 mol/L；总金属浓度m+ML=0.003 mol/L。
\item 配体束缚比例ML/CL=0.2324809832；核查l+HL+ML=0.012 mol/L等于0.01200。
\end{enumerate}
\noindent\textbf{答案摘要}\par
L\_free = 0.00837293473 mol/L；HL = 0.000837293473 mol/L；M\_free = 0.000210228201 mol/L；ML = 0.0027897718 mol/L；ligand\_bound\_fraction = 0.232480983 fraction\par
\noindent\textbf{模型前提：}
缓冲固定pH=9，只有HL、L−、M3+与一配体ML2+四种指定物种；忽略其它酸碱及配位形态。\par
\medskip\hrule\medskip
\Needspace{7\baselineskip}
\subsubsection*{UC-ADV5-I23 · 金属间交换氧化还原与有限电子容量的反演}
\noindent\textbf{题目}\par
假想一电子金属交换反应A(s)+B+(aq)⇌A+(aq)+B(s)，固相活度均为1，298.15K的平衡常数为K=[A+]/[B+]=2.300。初始B+物质的量0.05000mol、A+为0，纯固体A和B均充分过量，溶液体积保持1.000L。求平衡反应进度、两溶解金属量、A被氧化的电子mol、外电路理论电量及以初态反应商Q→0时电势发散说明为什么不能用一恒电势计算做功。\par
\medskip\noindent\textbf{分步解答}
\begin{enumerate}
\item 设反应进度ξ，溶液n(A+)=ξ、n(B+)=n0−ξ，物理区间0<ξ<n0。
\item 固态活度为1，反应商Q=ξ/(n0−ξ)，平衡要求Q=K=2.300。
\item 直接得ξ=K n0/(1+K)=0.03484848485mol，nB=0.01515151515mol。
\item 正向反应每前进一步发生一次一电子转移，n(e−)=ξ=0.03484848485mol。
\item 按F=96485.33212 C/mol，理论电量Qext=ξF=3362.367634C。
\item 初始A+为0时模型的浓度型Nernst电势因lnQ发散，真实体系需活度与成核约束；能量应积分E(ξ)dQ而不是固定EξF。
\item 数值逆核验：按独立化学计量、平衡和约束关系复算，三个无量纲检验残差分别为 2.220e-16、0.000e+0、0.000e+0；满足题定理想模型，但不代表专家审定。
\end{enumerate}
\noindent\textbf{答案摘要}\par
xi = 0.0348484848 mol；Aplus = 0.0348484848 mol；Bplus = 0.0151515152 mol；charge = 3362.36763 C\par
\noindent\textbf{模型前提：}
一电子假想金属交换，无离子迁移过电位与活度修正，明确不能以Q=0推断真实无穷电势。\par
\medskip\hrule\medskip
\Needspace{7\baselineskip}
\subsubsection*{UC-ADV5-I24 · 晶面弹性微应变对布拉格峰位的非线性移动}
\noindent\textbf{题目}\par
某立方晶体(110)面的一阶XRD反射用λ=0.1540nm X射线，初始Bragg角θ=29.30°（仪器记录2θ）。在与(110)面法向一致的弹性拉伸下，该面间距相对增加0.2000\%，其他晶面不考虑。求初始面间距、晶胞边长a、应变后峰位2θ、峰位差与线性小应变近似Δ(2θ)≈−2tanθ·ε（弧度）的相对误差。须用精确Bragg方程和线性展开独立复核。\par
\medskip\noindent\textbf{分步解答}
\begin{enumerate}
\item 初始Bragg条件2d sinθ=λ，求d=0.1573411546nm。
\item 立方晶体d110=a/√2，晶胞边长a=√2d=0.2225139948nm。
\item 拉伸后d′=d(1+ε)，保持同一X射线波长，θ′=arcsin[λ/(2d′)]。
\item 新2θ′=58.47168543度，峰位变化Δ2θ=-0.1283145697度，为负表示峰向低角度移动。
\item 小应变一阶式Δ(2θ)≈−2tanθ·ε（弧度）=-0.1286115676度。
\item 精确位移与一阶估计相比较，二者差来自arcsin的非线性及高阶应变项，不能混淆θ与2θ。
\item 数值逆核验：按独立化学计量、平衡和约束关系复算，三个无量纲检验残差分别为 0.000e+0、0.000e+0、0.000e+0；满足题定理想模型，但不代表专家审定。
\end{enumerate}
\noindent\textbf{答案摘要}\par
d = 0.157341155 nm；a = 0.222513995 nm；2theta\_new = 58.4716854 degree；shift = -0.12831457 degree\par
\noindent\textbf{模型前提：}
单晶法向应变，忽略仪器宽化、泊松应变、峰重叠及各向异性弹性效应。\par
\medskip\hrule\medskip
\Needspace{7\baselineskip}
\subsubsection*{UC-ADV5-I25 · 晶体场多重态的简并配分函数与磁响应反演}
\noindent\textbf{题目}\par
假想过渡金属配合物在300 K有两个能级：低能态E0=0、简并度g0=1，磁响应为1.5任意单位；高能态E1=0.1200 eV、简并度g1=5，磁响应为4.5同单位。分子之间独立、每种微态按Boltzmann加权，kB=1.380649×10\textasciicircum{}-23 J/K，1 eV=1.602176634×10\textasciicircum{}-19 J。求高自旋态占比、热平均磁响应、构型熵（每摩尔）与温度趋0和无限高时的极限；证明简并度不可遗漏。\par
\medskip\noindent\textbf{分步解答}
\begin{enumerate}
\item 配分函数Z=g0e\textasciicircum{}(−E0/kT)+g1e\textasciicircum{}(−E1/kT)=1+5e\textasciicircum{}(−Δ/kT)。
\item 高态占比p1=5e\textasciicircum{}(−Δ/kT)/Z=0.04598477196；低态占比为1−p1。
\item 热平均磁响应χ=1.5(1−p1)+4.5p1=1.637954316。
\item 按状态概率与内部简并度，摩尔构型熵S=−RΣpi ln(pi/gi)=2.166151092 J/(mol K)。
\item 低温下p1→0，χ→1.5、S→0；高温下p1→5/6，χ→(1.5+5×4.5)/6。
\item 若忽略g1=5，则得到错误的p1=e\textasciicircum{}(−Δ/kT)/(1+e\textasciicircum{}(−Δ/kT))，破坏高温简并度极限。
\item 数值逆核验：按独立化学计量、平衡和约束关系复算，三个无量纲检验残差分别为 0.000e+0、5.551e-17、0.000e+0；满足题定理想模型，但不代表专家审定。
\end{enumerate}
\noindent\textbf{答案摘要}\par
high\_frac = 0.045984772 fraction；chi = 1.63795432 arb；Sm = 2.16615109 J/(mol K)\par
\noindent\textbf{模型前提：}
理想两态配位场配分模型，不等价于真实磁化率的全部温度依赖；需要额外考虑Curie因子与交换。\par
\medskip\hrule\medskip
\Needspace{7\baselineskip}
\subsubsection*{UC-ADV5-I26 · 异价盐离子对形成与自由离子强度的闭合反演}
\noindent\textbf{题目}\par
在支持电解质充分屏蔽的理想水溶液里，1:1阳离子A+、阴离子B−的总浓度分别为0.004000、0.002000mol/L，其他惰性反离子用于电中性。只形成中性AB离子对，其形成常数K=[AB]/([A+][B−])=200 L/mol，体积恒定。求离子对、两种游离离子、B配对分数及这两类离子对离子强度I=1/2([A+]+[B−])的贡献，证明离子对形成导致自由载流离子总数下降。\par
\medskip\noindent\textbf{分步解答}
\begin{enumerate}
\item 设x=[AB]，则[A+]=0.004−x，[B−]=0.002−x，必须0<x<0.002。
\item 质量作用x=200(0.004−x)(0.002−x)，属于二次方程而非完全配对。
\item 将左端除以两种游离物浓度得到x/[(a−x)(b−x)]，在物理区间严格递增，故唯一根。
\item 求得x=0.000783009434 M，自由A+=0.003216990566 M、自由B−=0.001216990566 M。
\item B的配对比例x/0.002=0.391504717，自由两离子贡献离子强度为(0.003216990566+0.001216990566)/2=0.002216990566 M。
\item 与未配对的(a+b)/2比较，离子对形成使两种自由带电粒子合计减少2x，配对的中性AB本身不进入离子强度公式。
\item 数值逆核验：按独立化学计量、平衡和约束关系复算，三个无量纲检验残差分别为 0.000e+0、0.000e+0、0.000e+0；满足题定理想模型，但不代表专家审定。
\end{enumerate}
\noindent\textbf{答案摘要}\par
AB = 0.000783009434 mol/L；A\_free = 0.00321699057 mol/L；B\_free = 0.00121699057 mol/L；Ipartial = 0.00221699057 mol/L\par
\noindent\textbf{模型前提：}
仅对所列A、B计离子强度贡献，不把其余支持电解质忽略于总I；形成常数已是题定条件常数。\par
\medskip\hrule\medskip
\Needspace{7\baselineskip}
\subsubsection*{UC-ADV5-I27 · 固态碳酸盐热分解的动力学—产气容量反演}
\noindent\textbf{题目}\par
某假想碳酸盐晶体在恒温固定气氛下发生MCO3(s)→MO(s)+CO2(g)，有效转化分数遵循一级固相经验律dX/dt=k(1−X)，k=1.20×10\textasciicircum{}-3 s\textasciicircum{}-1。起始纯MCO3有0.05000mol，气体及时排出且不逆反应，600.0s后CO2按298.15K、101325Pa理想气体计量。求转化分数、剩余碳酸盐、生成氧化物及CO2的mol与体积，给出若初始夹带0.010mol不可分解杂质时的总残渣mol讨论。\par
\medskip\noindent\textbf{分步解答}
\begin{enumerate}
\item 微分方程dX/dt=k(1−X)分离变量并由X(0)=0得X(t)=1−e\textasciicircum{}(−kt)。
\item 600s转化X=0.513247744，残余碳酸盐n=(1−X)n0=0.0243376128mol。
\item 按MCO3→MO+CO2，一摩尔固体反应生成一摩尔MO与一摩尔CO2，均为0.0256623872mol。
\item 按PV=nRT，气体体积V=0.6278406627 L（R使用SI时先算m³）。
\item 初始0.010mol惰性杂质若不反应，最终固相物种mol总数为未反应碳酸盐+生成MO+惰性杂质=0.06mol。
\item 逆代入动学式得1−X=e\textasciicircum{}(−kt)，碳与氧质量账本须根据M原子数而非直接把不同物相mol相加当质量。
\item 数值逆核验：按独立化学计量、平衡和约束关系复算，三个无量纲检验残差分别为 0.000e+0、0.000e+0、0.000e+0；满足题定理想模型，但不代表专家审定。
\end{enumerate}
\noindent\textbf{答案摘要}\par
X = 0.513247744 fraction；unreacted = 0.0243376128 mol；CO2 = 0.0256623872 mol；Vgas = 0.627840663 L\par
\noindent\textbf{模型前提：}
经验一级固相分解、外部及时排除气体，不断面积或活化能，忽略反向碳酸化。\par
\medskip\hrule\medskip
\Needspace{7\baselineskip}
\subsubsection*{UC-ADV5-I28 · 混合价铁氧化物的电荷补偿与穆斯堡尔双谱线占比}
\noindent\textbf{题目}\par
某假想电荷补偿型铁氧化物每一化学式单位含2个Fe与3个氧，另有可占据的Li+填隙离子δ个。化学式Li\_δFe2O3，所有O为−2，Fe只含+2/+3，两Fe平均价由电中性决定。热重—化学滴定测得δ=0.4000（每式单位）；求Fe2+个数及Fe3+个数、Fe2+占铁原子比例、Fe平均价与理论磁响应强度比（假定谱线原子响应相同）；验证总电荷。\par
\medskip\noindent\textbf{分步解答}
\begin{enumerate}
\item 每化学式单位Li+填隙离子贡献δ正电，3个氧以O2−计总负电−6，因此Fe总正价必须为6−δ。
\item 设Fe(II)个数为x，Fe(III)个数为2−x，则铁总正价2x+3(2−x)=6−x。
\item 由电中性6−x+δ−6=0得到x=δ=0.4000个Fe(II)，Fe(III)=1.6000。
\item Fe(II)占全部铁原子分数x/2=0.2000，Fe(III)占0.8000；响应因子相同的谱峰面积比FeII/FeIII=0.4000/1.6000=0.2500。
\item 平均铁氧化态=(2×0.4+3×1.6)/2=2.8000；3个O2−负价总和−6，Li+与两类Fe的正价总和同为+6。
\item 若Li替代了晶格Fe位点或出现氧缺陷，该元素及位点总数会改变，必须重新列位点电中性，不能沿用上述x=δ。
\item 独立数值检验：(i) Fe(II)+Fe(III)=0.4+1.6=2；(ii) Li+与两种Fe正价合计0.4+2×0.4+3×1.6=6.0；(iii) 平均Fe价与谱学原子数分数给出2.8和0.2，互相吻合。
\end{enumerate}
\noindent\textbf{答案摘要}\par
FeII = 0.4 atoms/formula；FeIII = 1.6 atoms/formula；fractionFeII = 0.2 fraction；meanFe = 2.8 valence\par
\noindent\textbf{模型前提：}
题设Li填隙且Fe仅有II、III；真实LiδFe2O3的结构稳定性不予声称。\par
\medskip\hrule\medskip
\Needspace{7\baselineskip}
\subsubsection*{UC-ADV5-I29 · 双金属氧化价态反演与氧化剂电子当量审计}
\noindent\textbf{题目}\par
独立定量分析显示1.000L溶液中含0.01000mol Fe2+及0.02000mol Cr3+，无其它可氧化金属。施加强氧化剂使全部Fe2+→Fe3+、全部Cr3+→Cr6+（题设酸性可稳定），氧化剂每mol接受5mol电子，电流收集效率80.0\%。求需要的氧化剂物质的量、金属各放出电子总mol、有效电量、实际通过电量及氧化剂/金属电子当量闭合；讨论为何不能只数被氧化的金属原子个数。\par
\medskip\noindent\textbf{分步解答}
\begin{enumerate}
\item Fe(II)→Fe(III)每原子失1电子，电子mol为0.01000。
\item Cr(III)→Cr(VI)每原子失3电子，电子mol为3×0.02000=0.06000。
\item 总放出电子ne=0.07000mol，对应氧化剂每mol接受5电子，理论氧化剂nox=0.01400mol。
\item 有效法拉第电量Qeff=ne F=6753.973248C；80\%效率时通过电路的实际总电量Qtotal=Qeff/0.8=8442.466561C。
\item 按每mol氧化剂接受5mol电子，5nox=0.07000mol，必须与金属失电子和相等。
\item 氧化数变化而非被氧化的金属原子计数决定总电子转移；铬的实际含氧配体与氢离子需求不影响本题电子当量闭合。
\item 数值逆核验：按独立化学计量、平衡和约束关系复算，三个无量纲检验残差分别为 -1.983e-16、0.000e+0、0.000e+0；满足题定理想模型，但不代表专家审定。
\end{enumerate}
\noindent\textbf{答案摘要}\par
ne = 0.07 mol；oxidizer = 0.014 mol；Qeff = 6753.97325 C；Qtotal = 8442.46656 C\par
\noindent\textbf{模型前提：}
假想完全选择性氧化、氧化剂参与反应未受其它副反应消耗、80\%电流效率独立给定。\par
\medskip\hrule\medskip
\Needspace{7\baselineskip}
\subsubsection*{UC-ADV5-I30 · 酸碱缓冲条件下单配位金属的条件常数误用辨析}
\noindent\textbf{题目}\par
一假想二价金属M2+在pH5.000与单负配体L−发生唯一的1:1络合ML+；HL的pKa=4.000，游离配体酸碱平衡快速建立。金属分析浓度CM=0.02000mol/L、配体分析浓度CL=0.01000mol/L、对L−形成常数β=200 L/mol。计算自由L−与HL、自由M2+与ML+，求配体配位分数，并与错误地把CL当自由L−所得ML+对比差异及物理不可行性。\par
\medskip\noindent\textbf{分步解答}
\begin{enumerate}
\item 由pH−pKa=1可得HL/L−=10\textasciicircum{}(pKa−pH)=0.1。
\item 设l=[L−]，金属自由浓度m=CM/(1+βl)，络合ML=CMβl/(1+βl)。
\item 配体守恒CL=1.1l+ML，左端单调递增，所以唯一物理解存在。
\item 用二分求游离L−=0.00270610816 M，HL=0.000270610816 M。
\item 金属游离[M2+]=0.01297671898 M，络合[ML+]=0.007023281024 M，配体配位率[ML]/CL=0.7023281024。
\item 若错误令l=CL，则络合量为CMβCL/(1+βCL)=0.01333333333 M；与CL对比可见会消耗超过可用配体，因此违背配体守恒。
\item 数值逆核验：按独立化学计量、平衡和约束关系复算，三个无量纲检验残差分别为 1.735e-16、0.000e+0、2.220e-16；满足题定理想模型，但不代表专家审定。
\end{enumerate}
\noindent\textbf{答案摘要}\par
L = 0.00270610816 mol/L；HL = 0.000270610816 mol/L；Mfree = 0.012976719 mol/L；ML = 0.00702328102 mol/L\par
\noindent\textbf{模型前提：}
缓冲固定pH、只允许1:1物种且忽略离子活动系数；两个pKa描述仅有一个。\par
\medskip\hrule\medskip
\Needspace{7\baselineskip}
\subsubsection*{UC-ADV5-I31 · 氧缺失导致的CeO2−δ中Ce(III)比例与磁矩假设}
\noindent\textbf{题目}\par
某理想混合价氧化物CeO\_(2−δ)，δ=0.2000个氧原子/化学式单位。氧化铈中只允许Ce(IV)及Ce(III)，氧一律为O2−且其它缺陷不承担电荷。假定总Ce位点1个，Ce(III)的有效磁响应（不包含单位换算）为每CeIII原子2.540μB、CeIV为0。求CeIII比例、平均Ce价、每mol氧损失mol以及假设自旋完全平行时的化学式净磁矩；讨论无长程磁序时为何不能将净矩当真实测量值。\par
\medskip\noindent\textbf{分步解答}
\begin{enumerate}
\item 氧负电荷为−2(2−δ)=−(4−2δ)，所以Ce平均氧化态须为4−2δ。
\item 若CeIII分数为x，则平均价=3x+4(1−x)=4−x，因此x=2δ=0.4。
\item CeIV分数1−x=0.6，平均Ce价=3.6。
\item 每mol化学式相对于CeO2少δmol氧原子，等效氧气损失δ/2mol O2。
\item 若CeIII所有磁矩完全同向，化学式净磁矩在该**额外磁性假设**下为2.540x=1.016 μB。
\item 通常热涨落导致随机自旋时饱和磁矩、有效顺磁矩和零场净磁化必须区分，化学价态比例并不能唯一给出真实零场净磁矩。
\item 数值逆核验：按独立化学计量、平衡和约束关系复算，三个无量纲检验残差分别为 0.000e+0、0.000e+0、0.000e+0；满足题定理想模型，但不代表专家审定。
\end{enumerate}
\noindent\textbf{答案摘要}\par
CeIII\_fraction = 0.4 fraction；CeIV\_fraction = 0.6 fraction；mean\_ce = 3.6 valence；aligned\_mu = 1.016 mu\_B\par
\noindent\textbf{模型前提：}
理想固定Ce位且价态仅III、IV；磁性结果为明示饱和共线假设，不声称实验零场磁化。\par
\medskip\hrule\medskip
\Needspace{7\baselineskip}
\subsubsection*{UC-ADV5-I32 · 有限双齿配体逐级络合的物种分布与根唯一性}
\noindent\textbf{题目}\par
假想金属M2+总分析浓度0.02000M与一个单齿中性配体L总量0.01000M可形成ML2+和ML2\textasciicircum{}2+，逐级常数K1=300L/mol、K2=80L/mol，体积和pH固定，无其它物种。求自由配体、游离金属、两种络合物浓度、金属平均结合配体数、配体占用分数，证明完整配体守恒物理解唯一，讨论把二级络合物当作只消耗一个L的错误。\par
\medskip\noindent\textbf{分步解答}
\begin{enumerate}
\item 逐级配合关系ML=K1 l m，ML2=K1K2 l²m，令l=[L]、m=[M2+]。
\item 金属守恒CM=m(1+K1l+K1K2l²)，故m=CM/(1+K1l+K1K2l²)。
\item 配体守恒CL=l+ML+2ML2，代入单变量函数f(l)=l+CM(K1l+2K1K2l²)/(1+K1l+K1K2l²)。
\item f(l)对l≥0严格递增（配体平均占据数是正活动度的单调函数），f(0)=0且f(CL)>CL，唯一根。
\item 二分求自由l=0.001717449017M，游离m=0.01261013493M，ML=0.006497179152M，ML2=0.0008926859157M。
\item 平均每金属结合配体数(ML+2ML2)/CM=0.4141275492，配体占用比例(ML+2ML2)/CL=0.8282550983；两金属络合物共存。
\item 数值逆核验：按独立化学计量、平衡和约束关系复算，三个无量纲检验残差分别为 -1.735e-16、-1.735e-16、0.000e+0；满足题定理想模型，但不代表专家审定。
\end{enumerate}
\noindent\textbf{答案摘要}\par
Lfree = 0.00171744902 mol/L；Mfree = 0.0126101349 mol/L；ML = 0.00649717915 mol/L；ML2 = 0.000892685916 mol/L\par
\noindent\textbf{模型前提：}
假想逐级常数只允许单核ML与ML2，忽略聚合、配体酸碱及其它金属竞争。\par
\medskip\hrule\medskip
\subsection*{分析化学}
\addcontentsline{toc}{subsection}{第5辑 / 分析化学}
\Needspace{7\baselineskip}
\subsubsection*{UC-ADV5-A20 · 改变终体积的标准加入与竞争络合吸光响应联合反演}
\noindent\textbf{题目}\par
10.00mL未知金属试液原浓度c未知。分别向三份等体积原液中加入标准金属0、2.000、5.000mL（浓度0.01000M），体积仅由加标增加。显色前添加足量络合掩蔽剂，条件下金属的自由比例固定为1/(1+βL)，β=500 L/mol、自由L=0.002000M，只有游离金属对检测波长吸光，ε=7000 L mol\textasciicircum{}-1cm\textasciicircum{}-1、光程1.000cm，读数依次为5.6、10.5、15.4。求原液c，补偿稀释后再计算两次加标的理论吸光度、确定零加标曲线截距，说明忽视体积将产生什么偏差。\par
\medskip\noindent\textbf{分步解答}
\begin{enumerate}
\item 设未知原液浓度c，加入标准体积v（L）后金属物料为c V0+cs v，总体积V0+v。
\item 显色自由金属比例α=1/(1+βL)=1/2，因此A(v)=εα(cV0+csv)/(V0+v)。
\item v=0时A0=εαc，故c=A0/(εα)=0.0016 M。
\item 把求得c代入v=.002L，预测A2=10.5；v=.005L，预测A5=15.4。
\item 如构造校正坐标A(v)(V0+v)对v，得到直线εα(cV0+csv)，截距=0.056 吸光度·L，斜率=35 吸光度/L。
\item 如果直接将原始A对加标体积强迫线性拟合而不考虑最终体积改变，将引入非线性系统偏差；掩蔽剂比例也不可遗漏。
\item 最后将识别出的未知量代回相互独立的测量方程和物料约束，归一化残差为 0.000e+0、0.000e+0、0.000e+0；不以此代替现实仪器验证。
\end{enumerate}
\noindent\textbf{答案摘要}\par
c = 0.0016 mol/L；A2 = 10.5 abs；A5 = 15.4 abs；slope = 35 abs/L\par
\noindent\textbf{模型前提：}
自由L浓度由过量缓冲固定且加标不改变，只有游离金属吸光，背景扣除，标准加入对体积影响全部计入。\par
\medskip\hrule\medskip
\Needspace{7\baselineskip}
\subsubsection*{UC-ADV5-A21 · 有限蛋白结合位点荧光猝灭的二次反演}
\noindent\textbf{题目}\par
某荧光探针P与有限蛋白位点L发生1:1可逆结合PL，条件形成常数K=[PL]/([P][L])=240L/mol。总体积固定，探针总浓度0.002000M、蛋白可占位点总浓度0.003000M。只有自由P发光，PL在检测窗口中无发光，零蛋白条件下归一化信号为1，且探针的吸收与量子产率与浓度无关。求结合复合物、自由P/L、荧光比值、表观占据率，证明正根唯一以及为什么不能将游离蛋白浓度当作总蛋白。\par
\medskip\noindent\textbf{分步解答}
\begin{enumerate}
\item 设复合物浓度x，两个元素物料守恒分别给[P]=P0−x和[L]=L0−x。
\item 结合平衡要求x=K(P0−x)(L0−x)，物理区间0<x<min(P0,L0)。
\item 函数x/[(P0−x)(L0−x)]在物理区间单调从0至∞，所以唯一正根。
\item 解得复合物x=0.0007094535669 M，自由P=0.001290546433 M，自由L=0.002290546433 M。
\item 仅自由P发光，归一化相对荧光S/S0=[P]/P0=0.6452732165，位点占据率x/L0=0.2364845223。
\item 注意探针结合比例x/P0=0.3547267835与蛋白占据率x/L0不同；二者只在两种总浓度相等时才可能相等。
\item 最后将识别出的未知量代回相互独立的测量方程和物料约束，归一化残差为 2.220e-16、0.000e+0、0.000e+0；不以此代替现实仪器验证。
\end{enumerate}
\noindent\textbf{答案摘要}\par
PL = 0.000709453567 mol/L；freeP = 0.00129054643 mol/L；freeL = 0.00229054643 mol/L；fluorescence\_ratio = 0.645273217 fraction\par
\noindent\textbf{模型前提：}
理想单一位点、无协同、荧光仅自由探针、无内滤效应与复合物散射。\par
\medskip\hrule\medskip
\Needspace{7\baselineskip}
\subsubsection*{UC-ADV5-A22 · 顶空气液分配与有限水相回收的守恒反演}
\noindent\textbf{题目}\par
某水样中可挥发中性分析物初始浓度0.004000M、体积100.0mL；将样品放入密闭125mL容器，顶空气相体积25.0mL。达到298.15K平衡时，题设温度下气相浓度/水相浓度的无量纲分配系数Kcg/cw=3.000，两相体积固定、理想气体、没有瓶壁吸附。求平衡顶空气与水相浓度、转入顶空气的mol和占比、气相分压（kPa）；再说明仅凭顶空气信号反演原样浓度时必须校正哪些体积和分配参数。\par
\medskip\noindent\textbf{分步解答}
\begin{enumerate}
\item 起始分析物总mol数n0=c0 Vw=0.0004 mol，封闭容器不允许损失。
\item 分配平衡cg=Kcw，封闭物料守恒cwVw+cgVg=n0。
\item 消去cg得到cw=n0/(Vw+KVg)=0.002285714286 M，cg=0.006857142857 M。
\item 顶空气分析物mol为cgVg=0.0001714285714 mol，占总量的0.4285714286。
\item 理想气体分压p=cgRT，cg由mol/L换算为mol/m³时乘1000，Pa换算为kPa再除1000，两因子相消，故p=16.99856249 kPa。
\item 如仅使用顶空气响应cg，原始浓度c0=cg(Vw+KVg)/(KVw)，必须知道液相/气相体积与分配K，否则无法还原原水样。
\item 最后将识别出的未知量代回相互独立的测量方程和物料约束，归一化残差为 -1.355e-16、0.000e+0、0.000e+0；不以此代替现实仪器验证。
\end{enumerate}
\noindent\textbf{答案摘要}\par
cw = 0.00228571429 mol/L；cg = 0.00685714286 mol/L；gas\_fraction = 0.428571429 fraction；partial\_pressure = 16.9985625 kPa\par
\noindent\textbf{模型前提：}
Kcg/cw使用两相浓度同单位，气相理想，外部温度恒定、无反应、无吸附。\par
\medskip\hrule\medskip
\Needspace{7\baselineskip}
\subsubsection*{UC-ADV5-A23 · 双金属滴定终点的竞争络合选择性与残余游离金属}
\noindent\textbf{题目}\par
100.0mL溶液含两个不同金属M2+与N2+，分析浓度分别0.02000M、0.01500M；加入单齿络合剂L总浓度0.02500M（以最终混合体积计），只形成ML、NL，条件形成常数分别10\textasciicircum{}4与700 L/mol。求自由L、游离两金属、两络合物以及优先被络合金属的占用比例；分析当L总量小于两金属总量时为何不能认为二者均完全络合，并证明唯一物理解。\par
\medskip\noindent\textbf{分步解答}
\begin{enumerate}
\item 令l为游离络合剂浓度，单金属络合平衡给mA=CA/(1+K1l)、mB=CB/(1+K2l)。
\item 配位物种为[ML]=CA−mA，[NL]=CB−mB，各消耗1份配体。
\item 完整配体物料方程l+CAK1l/(1+K1l)+CBK2l/(1+K2l)=CL，在l∈[0,CL]有唯一正根。
\item 数值根l=0.0009445272208 M，游离金属M=0.001914741866 M、N=0.009029785355 M。
\item 络合量ML=0.01808525813 M，NL=0.005970214645 M；配位剂用于M的比例ML/(ML+NL)=0.7518147034。
\item 由于CL<CA+CB且平衡形成常数有限，完全络合二金属会违反配体守恒，同时必须保留游离金属。
\item 最后将识别出的未知量代回相互独立的测量方程和物料约束，归一化残差为 -2.776e-16、0.000e+0、-3.331e-16；不以此代替现实仪器验证。
\end{enumerate}
\noindent\textbf{答案摘要}\par
L = 0.000944527221 mol/L；Mfree = 0.00191474187 mol/L；Nfree = 0.00902978535 mol/L；ML = 0.0180852581 mol/L；NL = 0.00597021465 mol/L\par
\noindent\textbf{模型前提：}
pH由缓冲固定，只有两种1:1配合物，忽略沉淀、配位异构和其它物种。\par
\medskip\hrule\medskip
\Needspace{7\baselineskip}
\subsubsection*{UC-ADV5-A24 · 三通道ICP-OES谱线重叠的满秩矩阵消干扰}
\noindent\textbf{题目}\par
三种金属离子X、Y、Z在三个多元素ICP-OES波长的扣背景响应为S1=ε(cX+cY)，S2=ε(cX+0.5cY+cZ)，S3=ε(0.2cX+cY+cZ)，其中ε=15000 L mol\textasciicircum{}-1cm\textasciicircum{}-1且仪器光程等效已经标定。测得读数S1、S2、S3依次为22.5、23.25、15。求三个元素浓度和各通道的重叠贡献，计算矩阵行列式是否非零，并用一个未参与求解的线性组合S3−S2+0.8S1验证Y的响应。\par
\medskip\noindent\textbf{分步解答}
\begin{enumerate}
\item 将各浓度缩放为U=εcX、V=εcY、W=εcZ，三个通道方程分别为S1=U+V、S2=U+0.5V+W、S3=0.2U+V+W。
\item 构造S3−S2+0.8S1=(0.5+0.8)V=1.3V，求得V=7.5。
\item 再由S1=U+V得到U=15，第三未知W=S2−U−0.5V=4.5。
\item 换回金属浓度cX=U/ε=0.001 M、cY=0.0005 M、cZ=0.0003 M。
\item 系数矩阵经消元后给出非零比例1.3，因此3个波长线性独立；若响应矩阵行列式为0，则存在不可辨识的线性组合。
\item 三通道中由X、Y、Z分别贡献的强度可依照响应系数矩阵计算；验证0.8S1+S3−S2=9.75=1.3V。
\item 最后将识别出的未知量代回相互独立的测量方程和物料约束，归一化残差为 0.000e+0、0.000e+0、0.000e+0；不以此代替现实仪器验证。
\end{enumerate}
\noindent\textbf{答案摘要}\par
X = 0.001 mol/L；Y = 0.0005 mol/L；Z = 0.0003 mol/L\par
\noindent\textbf{模型前提：}
仪器响应已完成线性区校准、无基体改变和散射，三个波长的空白已扣除。\par
\medskip\hrule\medskip
\Needspace{7\baselineskip}
\subsubsection*{UC-ADV5-A25 · 双指数荧光寿命曲线的积分强度与时间门控反演}
\noindent\textbf{题目}\par
某理想脉冲激发荧光信号可表示I(t)=2000exp(−0.001t)+50000exp(−0.030t)，t单位ns，I为每ns光子计数率，仪器响应已经去卷积。两个寿命成分分别反映长寿命及短寿命荧光，忽略背景。计算t=10、30ns时瞬时强度、0至无穷的两成分积分光子数、以光子数加权平均的荧光寿命、10ns时间门控后长寿命部分的光子占比，说明瞬时振幅和总积分贡献不能混同。\par
\medskip\noindent\textbf{分步解答}
\begin{enumerate}
\item 信号是双指数叠加，给定两个寿命τ1=1/0.001=1000ns和τ2=1/0.030≈33.333ns。
\item 在t=10ns时I10=2000e\textasciicircum{}(−0.01)+50000e\textasciicircum{}(−0.30)=39021.0107，在t=30ns时I30=22269.37405。
\item 积分强度∫0∞Ai e\textasciicircum{}(−t/τi)dt=Ai τi，长寿命贡献2.0×10\textasciicircum{}6 photons、短寿命贡献约1.66667×10\textasciicircum{}6 photons。
\item 总积分光子数N=2000/0.001+50000/0.030=3666666.667，与初始振幅I0=52000 不同。
\item 以光子数加权τph=ΣAiτi²/ΣAiτi=560.6060606 ns。
\item 10ns后门控积分长寿命部分分数为(A1 τ1 e\textasciicircum{}(−10/τ1))/ΣAiτi e\textasciicircum{}(−10/τi)=0.6159330904；门控后组分权重不等于t=0振幅比例。
\item 最后将识别出的未知量代回相互独立的测量方程和物料约束，归一化残差为 0.000e+0、0.000e+0、0.000e+0；不以此代替现实仪器验证。
\end{enumerate}
\noindent\textbf{答案摘要}\par
I10 = 39021.0107 counts/ns；I30 = 22269.3741 counts/ns；total\_photons = 3666666.67 counts；mean\_tau = 560.606061 ns\par
\noindent\textbf{模型前提：}
探测响应去卷积且光子强度不受饱和，题给寿命为模型参数，没有统计噪声或背景。\par
\medskip\hrule\medskip
\Needspace{7\baselineskip}
\subsubsection*{UC-ADV5-A26 · 库仑水分测定中漂移空白与结晶水计量联合反演}
\noindent\textbf{题目}\par
某无水环境下库仑卡尔费休测定一个盐的结晶水，已知干盐样品0.2000g、总一价水反应电量换算成水量系数为每mol H2O需要2mol e−。有效测量积分电量Q=32.000C，其中同一时间窗口由独立空白试验估计电量漂移Qblank=1.500C，法拉第常数F=96485.33212 C/mol。假定试样中唯一可测挥发性成分为水，求扣空白的水mol与水质量、样品中水质量分数以及若把32.000C全计为水将造成的偏差。若结晶盐的无水式量M0=160.0g/mol，反推出平均每式单位水合数并说明它可能不是整数的含义。\par
\medskip\noindent\textbf{分步解答}
\begin{enumerate}
\item 有效净水电量QH2O=Q−Qblank=30.500 C，漂移不属于样品水。
\item 题设2e−对应1分子水，nH2O=QH2O/(2F)=0.0001580551123 mol。
\item 水质量mH2O=18.015 n=0.002847362847 g。
\item 按原试样0.2000g计算水质量分数w=0.01423681424。
\item 无水盐质量m0=0.2−mH2O=0.1971526372 g，盐式单位n0=m0/160；水合数n水/n盐=0.1282702495。
\item 若误扣空白则水量被系统高估Qblank/QH2O=0.04918032787；平均水合数非整数可能表示混合水合程度或测量/题设模型不完整。
\item 最后将识别出的未知量代回相互独立的测量方程和物料约束，归一化残差为 0.000e+0、0.000e+0、0.000e+0；不以此代替现实仪器验证。
\end{enumerate}
\noindent\textbf{答案摘要}\par
water\_mol = 0.000158055112 mol；water\_mass = 0.00284736285 g；water\_fraction = 0.0142368142 fraction；hydration = 0.12827025 mol/mol\par
\noindent\textbf{模型前提：}
理想库仑KF反应与标定每水2电子相符，背景电量已独立测定，无其它氧化还原物与挥发性干扰。\par
\medskip\hrule\medskip
\Needspace{7\baselineskip}
\subsubsection*{UC-ADV5-A27 · Van Deemter柱效、最优线速度与理论塔板数}
\noindent\textbf{题目}\par
某HPLC色谱柱长1.200m，实验得到Van Deemter经验系数A=0.001000mm、B=0.02000mm²/s、C=0.01000s（u单位mm/s），板高H(u)=A+B/u+Cu。要求严格最小化板高，求最佳线速度、最小板高及理论塔板数；线速度若偏离最优至两倍，求塔板数损失比例与半径条件，说明A项无法通过改变u消除。\par
\medskip\noindent\textbf{分步解答}
\begin{enumerate}
\item 对H(u)=A+B/u+Cu求导H′(u)=−B/u²+C，正驻点u*=√(B/C)=1.414213562mm/s。
\item 二阶导数H″(u)=2B/u³>0，所以该驻点是u>0上的唯一全局最小。
\item 最小板高Hmin=A+2√(BC)=0.02928427125mm。
\item 柱长1.2m=1200mm，理论塔板数N=L/Hmin=40977.62891。
\item 两倍最优线速度时H(2u*)=0.03635533906mm，塔板数N2=33007.53152。
\item 相对塔板损失1−N2/N*=0.1944987447；固定A反映与流速无关的多路径涡流扩散贡献。
\item 独立复算检验：将结果带回观测方程，分别恢复三项约束；计算残差为 0.000e+0、2.220e-16、0.000e+0。
\end{enumerate}
\noindent\textbf{答案摘要}\par
u\_star = 1.41421356 mm/s；Hmin = 0.0292842712 mm；Nmax = 40977.6289 plates；N\_double = 33007.5315 plates\par
\noindent\textbf{模型前提：}
经验Van Deemter模式只在指定柱和流速范围有效，不承诺实际超高速不存在压力限制及非理想传质。\par
\medskip\hrule\medskip
\Needspace{7\baselineskip}
\subsubsection*{UC-ADV5-A28 · 单元弱酸部分中和后的精确电荷守恒与pH}
\noindent\textbf{题目}\par
某恒容水溶液体积1.000L，弱酸HA的分析浓度C=0.002000M，Ka=1.000×10\textasciicircum{}-5；加入NaOH所提供的Na+分析浓度为0.001500M（不计体积变化），不含其他电解质。25°C时Kw=10\textasciicircum{}-14，假设所有酸、碱迅速平衡。要求由电荷守恒而非Henderson近似精确计算pH、HA与A−浓度和剩余自由H+/OH−浓度；证明物理解在h>0唯一并与忽略水自解离近似比较。\par
\medskip\noindent\textbf{分步解答}
\begin{enumerate}
\item 令h=[H+]，弱酸物料守恒与Ka关系给[A−]=C Ka/(Ka+h)、[HA]=C h/(Ka+h)。
\item 水自解离给[OH−]=Kw/h，电荷守恒h+[Na+]=[A−]+[OH−]。
\item 因此单变量方程f(h)=h+0.001500−C Ka/(Ka+h)−Kw/h=0。
\item 当h>0时f′(h)=1+C Ka/(Ka+h)²+Kw/h²>0，并且两端从负无穷到正无穷，故唯一正根。
\item 求得h=0.000003304055288M，pH=5.480952694，[A−]=0.001503301029M，[HA]=0.0004966989713M。
\item 水自解离提供[OH−]=3.026583737e-9M；验证电荷守恒差值为零，而直接使用pH=pKa+log10(na/nh)只在h及OH可忽略时近似成立。
\item 独立复算检验：将结果带回观测方程，分别恢复三项约束；计算残差为 -6.873e-17、0.000e+0、2.220e-16。
\end{enumerate}
\noindent\textbf{答案摘要}\par
pH = 5.48095269 pH；Aminus = 0.00150330103 mol/L；HA = 0.000496698971 mol/L；H = 3.30405529e-06 mol/L\par
\noindent\textbf{模型前提：}
理想浓度酸碱模型，NaOH引入钠已计入，没有CO2吸收、钠络合及活度修正。\par
\medskip\hrule\medskip
\Needspace{7\baselineskip}
\subsubsection*{UC-ADV5-A29 · 光谱响应与氧化还原电子当量的联立浓度解混}
\noindent\textbf{题目}\par
未知溶液只含两种可被定量氧化的分析物X、Y，起始浓度cX、cY未知。显色光谱响应A=2.000cX+0.7000cY，观测值A=0.034（浓度以mol/L计）；另一独立库仑扫描每mol X、Y分别转移1及2mol电子，测得每升的电子当量Q/F=0.05mol/L。求两个分析物浓度、若取20.00mL样品的完全氧化所需电量、光谱中X贡献比例以及模型的可辨识矩阵行列式。\par
\medskip\noindent\textbf{分步解答}
\begin{enumerate}
\item 未知量x,y满足光谱方程2x+0.7y=0.034。
\item 库仑氧化对每mol X提供1mol电子、Y提供2mol电子，故x+2y=0.05。
\item 两响应矩阵行列式D=2×2−0.7×1=3.3≠0，唯一解。
\item 由克拉默法则x=(2A−0.7Q)/D=0.01M、y=(2Q−A)/D=0.02M。
\item 取0.02000L样品时n\_e=(x+2y)0.02000mol，电量Qcell=F(x+2y)0.02000=96.48533212C。
\item 光谱中X响应比例2x/A=0.5882352941，不能以电子数比例代替吸光比例。
\item 独立复算检验：将结果带回观测方程，分别恢复三项约束；计算残差为 0.000e+0、0.000e+0、0.000e+0。
\end{enumerate}
\noindent\textbf{答案摘要}\par
X = 0.01 mol/L；Y = 0.02 mol/L；charge = 96.4853321 C；X\_signal\_frac = 0.588235294 fraction\par
\noindent\textbf{模型前提：}
两独立响应严格线性、光谱经过空白校正，两分析物不发生耦合氧化。\par
\medskip\hrule\medskip
\Needspace{7\baselineskip}
\subsubsection*{UC-ADV5-A30 · 大分子有限位点捕获的双重检测信号与竞争结合量}
\noindent\textbf{题目}\par
目标离子M与固相功能化捕获位点S的平衡为M+S⇌MS，条件形成常数350 L/mol，取同一参考体积浓度。初始总M=0.001000M、总位点S=0.001800M；无其它结合物种。发光探测只对游离M响应S\_opt=0.2000+900[M]，而独立重量分析对MS的物质的量响应。求游离M、剩余自由位点、捕获量、光学读数、捕获效率和若把位点浓度误当常量时的偏差方向。\par
\medskip\noindent\textbf{分步解答}
\begin{enumerate}
\item 设MS的浓度为x，则[M]=Mtot−x，[S]=Stot−x。
\item 质量作用x=K(Mtot−x)(Stot−x)；在物理区间0<x<min(Mtot,Stot)，函数x/[(Mtot−x)(Stot−x)]严格单调上升。
\item 求根得捕获浓度x=0.0003384275921M，自由M=0.0006615724079M。
\item 自由位点Stot−x=0.001461572408M，捕获效率x/Mtot=0.3384275921。
\item 光学信号Sopt=0.2+900[M]=0.7954151671，不应把总M直接代入光学信号。
\item 若误设自由位点[S]=Stot，等效增强结合能力，通常会高估x；严格结果必须同时满足两物料守恒。
\item 独立复算检验：将结果带回观测方程，分别恢复三项约束；计算残差为 -3.331e-16、0.000e+0、0.000e+0。
\end{enumerate}
\noindent\textbf{答案摘要}\par
captured = 0.000338427592 mol/L；Mfree = 0.000661572408 mol/L；Sfree = 0.00146157241 mol/L；optical = 0.795415167 signal\par
\noindent\textbf{模型前提：}
假想单一1:1功能位点，形成常数按标准参照浓度定义，不讨论捕获颗粒的非平衡吸附阻力。\par
\medskip\hrule\medskip
\Needspace{7\baselineskip}
\subsubsection*{UC-ADV5-A31 · 交叉灵敏传感器二通道消干扰与噪声条件数}
\noindent\textbf{题目}\par
两个电化学传感器对分析物X、Y的标定响应为S1=0.7000cX+0.2000cY与S2=0.3000cX+0.8000cY，实测S1=0.0011、S2=0.0019，浓度以mol/L计。求cX、cY、响应矩阵的行列式；若仅第二通道有+1.000×10\textasciicircum{}-4的偏移，求反演两浓度的确定性变化，并解释为什么通道响应向量越共线，噪声越容易放大。\par
\medskip\noindent\textbf{分步解答}
\begin{enumerate}
\item 写成线性矩阵M c=S，其中M=[[0.7,0.2],[0.3,0.8]]。
\item 其行列式D=0.7×0.8−0.2×0.3=0.5，不为零，故唯一反演。
\item 克拉默法则cX=(0.8S1−0.2S2)/D=0.001M。
\item cY=(0.7S2−0.3S1)/D=0.002M。
\item 当只有S2出现偏移δ时，ΔcX=−0.2δ/D=-0.00004M，ΔcY=0.7δ/D=0.00014M。
\item 逆问题的误差放大随1/|D|增大，两个通道完全共线会使D=0，根本无法从两个观测恢复两个未知量。
\item 独立复算检验：将结果带回观测方程，分别恢复三项约束；计算残差为 3.943e-16、3.424e-16、4.337e-16。
\end{enumerate}
\noindent\textbf{答案摘要}\par
X = 0.001 mol/L；Y = 0.002 mol/L；det = 0.5 dimensionless；shiftX = -4e-05 mol/L\par
\noindent\textbf{模型前提：}
传感器线性且背景扣除，讨论的是确定性偏移而非统计标准差估计。\par
\medskip\hrule\medskip
\Needspace{7\baselineskip}
\subsubsection*{UC-ADV5-A32 · 饱和响应型电化学适配体传感器的逆标定与动态范围}
\noindent\textbf{题目}\par
某适配体电化学传感器的归一化信号由单站位点占据控制，S(c)=Kc/(1+Kc)，形成常数K=4.000×10\textasciicircum{}4 L/mol，测量条件保持固定。读数S=0.9523809524。求待测浓度、dS/dc局部灵敏度、采用线性低浓度近似S≈Kc反算时的高估/低估倍数，以及如果S的绝对偏移为0.0100时得到的浓度误差线性近似；证明接近饱和时反算不稳定的数学原因。\par
\medskip\noindent\textbf{分步解答}
\begin{enumerate}
\item 饱和响应S=Kc/(1+Kc)单调从0上升至1，反函数为c=S/[K(1−S)]。
\item 将观测S=0.9523809524代入，c=0.0005mol/L。
\item 微分dS/dc=K/(1+Kc)²=90.70294785 L/mol。
\item 如果误用低浓度线性式c\_lin=S/K，得到0.00002380952381M，与真实浓度之比为0.04761904762，严重低估。
\item 反函数导数dc/dS=1/[K(1−S)²]=0.011025 M/信号；信号偏移0.0100时一阶浓度偏移约0.00011025M。
\item 当S→1，dc/dS→∞；因此越接近饱和，即使相同绝对信号误差也会带来更大的反算浓度误差。
\item 独立复算检验：将结果带回观测方程，分别恢复三项约束；计算残差为 0.000e+0、-1.084e-15、0.000e+0。
\end{enumerate}
\noindent\textbf{答案摘要}\par
c = 0.0005 mol/L；sensitivity = 90.7029478 L/mol；linear\_estimate = 2.38095238e-05 mol/L；error\_delta = 0.00011025 mol/L\par
\noindent\textbf{模型前提：}
仅理想独立1:1单站吸附，响应无迟滞、无交叉物和仪器漂移。\par
\medskip\hrule\medskip
\subsection*{有机化学}
\addcontentsline{toc}{subsection}{第5辑 / 有机化学}
\Needspace{7\baselineskip}
\subsubsection*{UC-ADV5-O19 · 有限底物分子间可逆缩合的反应进度与最大转化边界}
\noindent\textbf{题目}\par
假想有机缩合反应A+B⇌C严格符合1:1:1计量且不生成小分子（如可逆环加成），在1.000L恒温理想溶液中初始A=0.06000mol、B=0.03000mol、C=0，反应平衡常数K=[C]/([A][B])=8.000 L/mol。求平衡反应进度、各物种浓度、限制底物B的转化率，证明物理解唯一，并判断把K当成纯摩尔比K=[C]/([A][B])但忽略体积会造成何种单位问题。\par
\medskip\noindent\textbf{分步解答}
\begin{enumerate}
\item 设每升反应进度x，恒定体积下[A]=0.060−x、[B]=0.030−x、[C]=x。
\item 理想平衡给x=8(0.060−x)(0.030−x)，物理区间0<x<0.030。
\item 函数x/[(A0−x)(B0−x)]在该区间由0严格递增至∞，故唯一正根。
\item 求得x=0.008726268674 M，[A]=0.05127373133 M、[B]=0.02127373133 M、[C]=0.008726268674 M。
\item 限制底物B的转化率x/0.030=0.2908756225，A转化率x/0.060=0.1454378112；二者不相等。
\item 若改写为物质的量，平衡常数必须按浓度标准态引入体积因子，只有题设体积1L时数值恰巧相等，不能忽略K的单位规范。
\item 核算：将物种量、相对速率或选择性反带原始独立方程，各约束归一化残差为 -5.551e-16、0.000e+0、0.000e+0。
\end{enumerate}
\noindent\textbf{答案摘要}\par
extent = 0.00872626867 mol/L；A = 0.0512737313 mol/L；B = 0.0212737313 mol/L；Bconversion = 0.290875622 fraction\par
\noindent\textbf{模型前提：}
只允许A+B↔C，溶液理想、无其它计量产物，K是明示浓度型平衡常数。\par
\medskip\hrule\medskip
\Needspace{7\baselineskip}
\subsubsection*{UC-ADV5-O20 · 串联—并联有机反应网络的中间体浓度与瞬时选择性}
\noindent\textbf{题目}\par
假想有机网络A→B(k1=0.180min\textasciicircum{}-1)、B→C(k2=0.0600min\textasciicircum{}-1)、A→D(k3=0.0300min\textasciicircum{}-1)严格按不可逆一级反应进行，初始纯A=1.000mol，恒温恒容。求5.000min时A、B、C、D各自mol、B相对于已消耗A的选择性、B瞬时净生成速率；给出一般解析式并用总物料守恒验证，解释为何不能套用只有A→B→C的B最优分离时间公式。\par
\medskip\noindent\textbf{分步解答}
\begin{enumerate}
\item A的消耗方程dA/dt=−(k1+k3)A，所以A(t)=e\textasciicircum{}(−(k1+k3)t)。
\item B的方程dB/dt=k1A−k2B，积分因子法得B(t)=k1[e\textasciicircum{}(−k2t)−e\textasciicircum{}(−(k1+k3)t)]/(k1+k3−k2)。
\item 副产物D(t)=k3[1−e\textasciicircum{}(−(k1+k3)t)]/(k1+k3)，C=1−A−B−D。
\item 5min时A=0.3499377491、B=0.4690565659、D=0.09286603584、C=0.08813964916 mol。
\item 相对于已消耗A的B选择性S\_B=B/(1−A)=0.7215563821，B净生成速率dB/dt=0.03484540089 mol/min。
\item 并联A→D改变A衰减率为k1+k3，因此若直接使用无旁路串联公式将同时破坏A、B与总物料的核验。
\item 核算：将物种量、相对速率或选择性反带原始独立方程，各约束归一化残差为 0.000e+0、0.000e+0、0.000e+0。
\end{enumerate}
\noindent\textbf{答案摘要}\par
A = 0.349937749 mol；B = 0.469056566 mol；C = 0.0881396492 mol；D = 0.0928660358 mol\par
\noindent\textbf{模型前提：}
全网不可逆一级、恒体积，所有物种均不参与其他消耗路径；数值按1mol初始A计。\par
\medskip\hrule\medskip
\Needspace{7\baselineskip}
\subsubsection*{UC-ADV5-O21 · Curtin–Hammett快速构象平衡下两产物通道选择性}
\noindent\textbf{题目}\par
假想底物A与B是快速互变构象，任意时刻[B]/[A]=Kconf=1.800，两构象分别不可逆生成产物P与Q，微观一级常数kA→P=0.0400min\textasciicircum{}-1、kB→Q=0.1500min\textasciicircum{}-1；构象交换远快于两反应，整个体系初始底物1.000mol。求整体底物有效衰减常数、10min时剩余量、P与Q累积产物、产物Q/P比例，并严格讨论为何产物比例不是单纯构象比例Kconf。\par
\medskip\noindent\textbf{分步解答}
\begin{enumerate}
\item 快速构象平衡时A分数fA=1/(1+Kconf)=0.3571428571，B分数fB=Kconf/(1+Kconf)=0.6428571429。
\item 总底物S=A+B，dS/dt=−(fA kA+fB kB)S，故kobs=(kA+Kconf kB)/(1+Kconf)=0.1107142857 min\textasciicircum{}-1。
\item 剩余S(t)=exp(−kobs t)，10min时S=0.3305019045mol。
\item 产物P的瞬时生成速率=fA kA S，积分得P=[fA kA/kobs](1−S)=0.08638685103mol。
\item 同理Q=[fB kB/kobs](1−S)=0.5831112444mol，Q/P=Kconf kB/kA=6.75。
\item 只有构象平衡快于不可逆反应时结果成立；Q/P由构象概率与反应势垒两者共同决定，而不等于Kconf本身。
\item 核算：将物种量、相对速率或选择性反带原始独立方程，各约束归一化残差为 0.000e+0、8.882e-16、0.000e+0。
\end{enumerate}
\noindent\textbf{答案摘要}\par
kobs = 0.110714286 min\textasciicircum{}-1；surviving = 0.330501905 mol；P = 0.086386851 mol；Q = 0.583111244 mol\par
\noindent\textbf{模型前提：}
构象快速互变，产物不回转且不存在彼此转化；仅两种构象贡献两条不可逆通路。\par
\medskip\hrule\medskip
\Needspace{7\baselineskip}
\subsubsection*{UC-ADV5-O22 · 连续有机保护—脱保护两步网络的中间体最佳截断}
\noindent\textbf{题目}\par
假想受保护有机底物A在连续不可逆一级条件下发生A→B(k1=0.200h\textasciicircum{}-1)，B→C(k2=0.0800h\textasciicircum{}-1)，B是目标单保护中间体、C为过度脱保护副产物。起始纯A=1.000mol，在4.000h测定。求A、B、C的量、4h B收率、B收率最大对应的唯一停止时间与峰值，证明极值条件为中间体瞬时生成与消耗速率相等。\par
\medskip\noindent\textbf{分步解答}
\begin{enumerate}
\item 由一级方程得到A(t)=e\textasciicircum{}(−k1t)，B(t)=k1[e\textasciicircum{}(−k2t)−e\textasciicircum{}(−k1t)]/(k1−k2)。
\item 全物料守恒给C(t)=1−A(t)−B(t)。
\item 4h时A=0.4493289641、B=0.4613667883、C=0.08930424762mol。
\item 中间体导数B′(t)=k1A−k2B，最大时刻满足k1A=k2B。
\item 解析化简求得t*=ln(k1/k2)/(k1−k2)=7.635756099h，唯一极大值因为B(0)=B(∞)=0。
\item 将t*代入B(t)得最大收率=0.5428835233，不可用长期转化率为目标B收率代替最优截断。
\item 核算：将物种量、相对速率或选择性反带原始独立方程，各约束归一化残差为 0.000e+0、6.939e-18、0.000e+0。
\end{enumerate}
\noindent\textbf{答案摘要}\par
A4 = 0.449328964 mol；B4 = 0.461366788 mol；C4 = 0.0893042476 mol；optimal\_h = 7.6357561 h\par
\noindent\textbf{模型前提：}
不可逆两级一级动力学，不存在旁路副反应，B能够在指定时刻瞬时分离。\par
\medskip\hrule\medskip
\Needspace{7\baselineskip}
\subsubsection*{UC-ADV5-O23 · 催化中间体有限占有下的饱和动力学与周转数}
\noindent\textbf{题目}\par
假想均相有机催化循环有空闲催化剂Cat与底物结合态CatS两种主要静态物种，结合条件常数K=250 L/mol（非0.250），催化剂总浓度0.01000M，底物自由浓度由储液维持0.05000M，CatS→Cat+P的产物形成速率常数kcat=0.400s\textasciicircum{}-1且配合平衡比周转快。求CatS、空闲Cat、瞬时产物生成速率、每催化剂周转频率TOF与1h内的形式周转数TON；说明为什么TON在底物被固定供应时可以大于初始底物/催化剂比例。\par
\medskip\noindent\textbf{分步解答}
\begin{enumerate}
\item 准平衡条件下[CatS]/[Cat]=K[S]=250×0.05000=12.5。
\item 催化剂守恒Ctot=[Cat]+[CatS]，故结合态分数θ=KS/(1+KS)=0.9259259259。
\item 求得[CatS]=Ctotθ=0.009259259259M、空闲Cat=0.0007407407407M。
\item 产物生成速率v=kcat[CatS]=0.003703703704M/s。
\item 周转频率TOF=v/Ctot=kcatθ=0.3703703704s\textasciicircum{}-1。
\item 若底物由储液维持且催化剂不失活，1小时形式TON=TOF×3600=1333.333333；催化剂可循环，因此TON不受初始反应池自由底物量限制。
\item 核算：将物种量、相对速率或选择性反带原始独立方程，各约束归一化残差为 0.000e+0、0.000e+0、0.000e+0。
\end{enumerate}
\noindent\textbf{答案摘要}\par
catS = 0.00925925926 mol/L；cat = 0.000740740741 mol/L；rate = 0.0037037037 mol/(L s)；TON = 1333.33333 turnovers\par
\noindent\textbf{模型前提：}
CatS与Cat快速准平衡，底物自由浓度由大储液钳定，假设产品不抑制催化剂。\par
\medskip\hrule\medskip
\Needspace{7\baselineskip}
\subsubsection*{UC-ADV5-O24 · 最概然缩聚链长分布的数均—重均聚合度与分散系数}
\noindent\textbf{题目}\par
理想A-A/B-B等当量、完全线性的逐步缩聚在官能团转化率p=0.9400时形成Flory最概然分布，链长为n的分子数分数Pn=(1−p)p\textasciicircum{}(n−1)（n≥1），无环化与分支。原始单体总数以10000mol为单位统计。计算数均聚合度DPn、重均聚合度DPw、分散系数Đ、链长n=2的分子数分数、聚合度≥10的分子数分数以及全部分子簇数量，并从几何级数独立推导两个均值。\par
\medskip\noindent\textbf{分步解答}
\begin{enumerate}
\item 数分布归一化Σ\_\{n≥1\}(1−p)p\textasciicircum{}(n−1)=1。
\item 几何级数给Σn Pn=1/(1−p)，即DPn=16.66666667。
\item 重均DPw=Σn²Pn/ΣnPn=(1+p)/(1−p)=32.33333333。
\item 分散系数Đ=DPw/DPn=1+p=1.94，随转化率增加而趋近2。
\item 链长2的数分数P2=(1−p)p=0.0564，链长至少10的数分数Σ\_\{n≥10\}Pn=p\textasciicircum{}9=0.5729948022。
\item 若最初10000mol单体没有环化，最终链簇的总mol数为10000/DPn=600mol；这一数不是聚合物总质量。
\item 核算：将物种量、相对速率或选择性反带原始独立方程，各约束归一化残差为 0.000e+0、2.220e-16、0.000e+0。
\end{enumerate}
\noindent\textbf{答案摘要}\par
DPn = 16.6666667 repeat units；DPw = 32.3333333 repeat units；dispersity = 1.94 dimensionless；P2 = 0.0564 fraction\par
\noindent\textbf{模型前提：}
理想Flory最概然链长分布，全部链无环化和交联，忽略端基质量对分子质量分布的影响。\par
\medskip\hrule\medskip
\Needspace{7\baselineskip}
\subsubsection*{UC-ADV5-O25 · 末端模型共聚链的Markov序列矩阵与稳态二联体频率}
\noindent\textbf{题目}\par
假想两单体A/B自由基共聚，进料摩尔分数fA=0.6000、fB=0.4000，反应活性比rA=2.000、rB=0.5000；进料在时间尺度上恒定，末端模型仅由链端种类决定下一单体。求条件跃迁矩阵P(A|A)、P(B|A)、P(A|B)、P(B|B)，求长链稳态A末端分数以及AA/AB/BA/BB四种有向二联体概率；验证马尔可夫平稳条件和AB=BA。\par
\medskip\noindent\textbf{分步解答}
\begin{enumerate}
\item 末端为A时竞争速率比kAA[A]/kAB[B]=rA fA/fB，因此PAA=rA fA/(rA fA+fB)=0.75。
\item 同一行互补PAB=1−PAA=0.25。
\item 末端B时PBB=rB fB/(rB fB+fA)=0.25，PBA=1−PBB=0.75。
\item 设长链末端为A的稳态分数πA，满足πA=πA PAA+(1−πA)PBA，解得πA=PBA/(PBA+PAB)=0.75。
\item 有向二联体概率为AA=πA PAA=0.5625、AB=πA PAB=0.1875、BA=(1−πA)PBA=0.1875、BB=(1−πA)PBB=0.0625。
\item 总概率为1且平稳二态链满足AB=BA，不等于按进料独立抽样时的fA fB；因此反应活性比包含链序列信息。
\item 复核：按原始反应守恒、速率分布或跃迁矩阵还原答案，归一化误差为0.000e+0、8.327e-17、-1.110e-16。
\end{enumerate}
\noindent\textbf{答案摘要}\par
A\_end\_fraction = 0.75 fraction；AA = 0.5625 fraction；AB = 0.1875 fraction；BB = 0.0625 fraction\par
\noindent\textbf{模型前提：}
进料恒定、末端模型和稳态长链极限成立，忽略倒数第二单体效应、链终止偏好及浓度漂移。\par
\medskip\hrule\medskip
\Needspace{7\baselineskip}
\subsubsection*{UC-ADV5-O26 · 动态动力学拆分中消旋与两路产物生成的耦合矩阵解}
\noindent\textbf{题目}\par
起始外消旋体R、S各0.5000mol。R和S分别不可逆以kR=0.200h\textasciicircum{}-1、kS=0.0500h\textasciicircum{}-1生成PR、PS，同时未反应底物双向消旋R⇌S，两方向速率常数均为krac=0.0800h\textasciicircum{}-1。体系闭合恒温。利用2×2线性矩阵指数求6h时剩余R、S、总产物mol及底物ee，写出不将消旋误算作净消耗的微分方程、特征值及总物料守恒。\par
\medskip\noindent\textbf{分步解答}
\begin{enumerate}
\item 未反应底物微分方程dR/dt=−(kR+krac)R+krac S，dS/dt=krac R−(kS+krac)S。
\item 矩阵M的迹为−(kR+kS+2krac)，记b=(kR+kS+2krac)/2=0.205。
\item 其本征值为−b±d，d=sqrt[(kR−kS)²+4krac²]/2=0.109658561；矩阵指数e\textasciicircum{}(Mt)=e\textasciicircum{}(−bt)[I cosh(dt)+(M+bI)sinh(dt)/d]。
\item 初始R0=S0=0.5000，代入得到R6=0.1836450621mol，S6=0.3248742725mol。
\item 未反应底物总量=0.5085193346mol，故总产物mol=0.4914806654，剩余底物ee=(R−S)/(R+S)=-0.2777263338。
\item 消旋两条反应改变R与S分布但互相抵消总底物数导数；产品总生成速率kR R+kS S，故总A+产物=1。
\item 复核：按原始反应守恒、速率分布或跃迁矩阵还原答案，归一化误差为0.000e+0、0.000e+0、-1.388e-17。
\end{enumerate}
\noindent\textbf{答案摘要}\par
R = 0.183645062 mol；S = 0.324874273 mol；product\_total = 0.491480665 mol；ee = -0.277726334 fraction\par
\noindent\textbf{模型前提：}
线性不可逆消耗加双向一阶消旋，不含产物回转、非对称催化剂失活或底物相互转化以外副反应。\par
\medskip\hrule\medskip
\Needspace{7\baselineskip}
\subsubsection*{UC-ADV5-O27 · 静态—动态双重荧光猝灭的寿命与强度分离}
\noindent\textbf{题目}\par
假想有机荧光探针初态寿命τ0=5.000ns，加入猝灭剂Q浓度0.08000M。动态双分子猝灭速率常数kq=4.000×10\textasciicircum{}9 L mol\textasciicircum{}-1s\textasciicircum{}-1，另有独立静态不发光复合效应Ks=6.000L/mol。发光寿命只受动态猝灭影响，荧光强度同时受两效应。求新寿命、寿命比、强度比、寿命与强度的Stern–Volmer增强因子及由同时测得两种比值反演静态常数的表达式，说明单独强度不能分离两机制。\par
\medskip\noindent\textbf{分步解答}
\begin{enumerate}
\item 激发态衰减速率为1/τ0+kq[Q]，故τ=1/(1/τ0+kqQ)=1.923076923ns。
\item 寿命比τ/τ0=1/(1+kqτ0Q)=0.3846153846，寿命Stern–Volmer因子τ0/τ=2.6。
\item 静态复合减少初始可发光探针占比，强度比I/I0=(τ/τ0)/(1+KsQ)=0.2598752599。
\item 强度Stern–Volmer因子I0/I=(1+kqτ0Q)(1+KsQ)=3.848。
\item 若分别测得寿命与强度，则Ks=[(I0/I)/(τ0/τ)−1]/Q=6 L/mol。
\item 单一强度数据只能确定两个效应因子的乘积，无法区分静态配位与动态激发态碰撞猝灭。
\item 复核：按原始反应守恒、速率分布或跃迁矩阵还原答案，归一化误差为0.000e+0、0.000e+0、-4.441e-16。
\end{enumerate}
\noindent\textbf{答案摘要}\par
tau = 1.92307692 ns；lifetime\_ratio = 0.384615385 fraction；intensity\_ratio = 0.25987526 fraction；Kst = 6 L/mol\par
\noindent\textbf{模型前提：}
动态猝灭满足简单Stern–Volmer，静态猝灭物种完全暗态且不改变剩余自由探针的τ0。\par
\medskip\hrule\medskip
\Needspace{7\baselineskip}
\subsubsection*{UC-ADV5-O28 · 有机底物与催化剂预配位竞争抑制的周转响应}
\noindent\textbf{题目}\par
某均相有机催化剂Cat总浓度0.008000M，与底物S（由大储液维持0.03000M）形成唯一反应活性络合物CatS，结合常数K\_S=200L/mol；另有抑制剂I自由浓度0.02000M，形成不活性CatI，K\_I=80.0L/mol，两结合相互排斥且都迅速达到准平衡。仅CatS以kcat=0.500s\textasciicircum{}-1形成P并恢复Cat。求自由Cat、CatS、CatI、产品生成初速及相对无抑制剂时速率保留比例。\par
\medskip\noindent\textbf{分步解答}
\begin{enumerate}
\item 设自由Cat浓度c，由准平衡CatS=KS[S]c=6c、CatI=KI[I]c=1.6c。
\item 催化剂守恒Ctot=c(1+6+1.6)=8.6c，因此自由c=0.0009302325581M。
\item 活性络合CatS=6c=0.005581395349M、抑制态CatI=1.6c=0.001488372093M。
\item 产品形成初速v=kcat[CatS]=0.002790697674M/s。
\item 无抑制剂时CatS分数=6/(1+6)，v0=0.003428571429M/s。
\item 速率保留比v/v0=(1+6)/(1+6+1.6)=0.8139534884，说明高底物浓度能通过竞争结合缓解抑制。
\item 复核：按原始反应守恒、速率分布或跃迁矩阵还原答案，归一化误差为0.000e+0、0.000e+0、0.000e+0。
\end{enumerate}
\noindent\textbf{答案摘要}\par
cat\_free = 0.000930232558 mol/L；catS = 0.00558139535 mol/L；catI = 0.00148837209 mol/L；rate = 0.00279069767 mol/(L s)\par
\noindent\textbf{模型前提：}
竞争配位与配体浓度准稳态钳定，产物无抑制，无协同性，催化剂不会失活。\par
\medskip\hrule\medskip
\Needspace{7\baselineskip}
\subsubsection*{UC-ADV5-O29 · 多级有机前药酶促活化的有限时间疗效组分反演}
\noindent\textbf{题目}\par
假想有机前药P在血浆模拟体外反应体系中依次不可逆一级转为活性A再转为无活性代谢物M，速率常数k1=0.01200min\textasciicircum{}-1、k2=0.005000min\textasciicircum{}-1。初始P浓度0.02000M、A及M均为零，体系恒体积；120.0min后取样。求三种物质的浓度、已消耗P中转为仍有效A的比例、A最大时刻以及该时刻浓度，证明A(t)的解析解，并注明该模型不能直接推断人体疗效或药代分布。\par
\medskip\noindent\textbf{分步解答}
\begin{enumerate}
\item P(t)=P0 exp(−k1 t)，得120min时P=0.004738555174M。
\item 活性A满足A′=k1P−k2A，积分因子解A=P0 k1[e\textasciicircum{}(−k2t)−e\textasciicircum{}(−k1t)]/(k1−k2)。
\item 120min时A=0.01069316151M，失活代谢产物M=P0−P−A=0.004568283315M。
\item 已消耗P对应的活性A比例A/(P0−P)=0.7006650833。
\item A瞬时最大满足k1P=k2A，化简t*=ln(k1/k2)/(k1−k2)=125.0669625min。
\item 峰值浓度A*=0.01070164493M；实际血浆蛋白结合、非均相代谢与体内清除均不在此题范围。
\item 复核：按原始反应守恒、速率分布或跃迁矩阵还原答案，归一化误差为0.000e+0、0.000e+0、0.000e+0。
\end{enumerate}
\noindent\textbf{答案摘要}\par
P120 = 0.00473855517 mol/L；A120 = 0.0106931615 mol/L；M120 = 0.00456828332 mol/L；t\_peak = 125.066962 min\par
\noindent\textbf{模型前提：}
纯体外假想一级连续化学反应，不是医学建议或真实药物疗效预报。\par
\medskip\hrule\medskip
\Needspace{7\baselineskip}
\subsubsection*{UC-ADV5-O30 · 加热有机裂解反应的温度跳变与活化焓敏感性}
\noindent\textbf{题目}\par
某假想单分子有机键裂解A→P服从Arrhenius动力学k(T)=kref exp[−Ea/R(1/T−1/Tref)]，已知Tref=310.0K时kref=0.001000s\textasciicircum{}-1、活化能Ea=45.00kJ/mol。将初始1.000mol纯A迅速升至300.0K保持300s，假设升温瞬间零持续时间，裂解为不可逆一级。求300K速率常数、300s转化率、残留底物、由速率温度灵敏度dlnk/dT=Ea/(RT²)得到每K相对变化率，以及如果把T误用摄氏度代入会破坏哪一步数学自洽性。\par
\medskip\noindent\textbf{分步解答}
\begin{enumerate}
\item 由参考温度310K的速率常数，Arrhenius式直接给k(300)=0.0005588003595s\textasciicircum{}-1。
\item 不可逆一级反应在恒温300K时A(t)=A0 e\textasciicircum{}(−k t)，300s残留=0.8456581257mol。
\item 转化分数X=1−e\textasciicircum{}(−kt)=0.1543418743，生成产物mol与转化分数相同。
\item 在温度300K，dlnk/dT=Ea/(RT²)=0.06013617752 K\textasciicircum{}-1，体现速率对温度的局部灵敏度。
\item 若1K小温差，线性相对速率变化约(dlnk/dT)ΔT，但不能用该导数跨大温差替代完整指数。
\item Arrhenius推导依赖绝对热力学温度；直接用摄氏度会造成错误的1/T、错误的无量纲指数与物理行为。
\item 复核：按原始反应守恒、速率分布或跃迁矩阵还原答案，归一化误差为0.000e+0、0.000e+0、0.000e+0。
\end{enumerate}
\noindent\textbf{答案摘要}\par
k300 = 0.000558800359 s\textasciicircum{}-1；conversion = 0.154341874 fraction；remaining = 0.845658126 mol；sensitivity = 0.0601361775 K\textasciicircum{}-1\par
\noindent\textbf{模型前提：}
纯一级热裂解，活化能恒定、瞬时温度跳变，不考虑热容、反应放热和压力反馈。\par
\medskip\hrule\medskip
\subsection*{物理化学}
\addcontentsline{toc}{subsection}{第5辑 / 物理化学}
\Needspace{7\baselineskip}
\subsubsection*{UC-ADV5-P19 · 正则溶液的自旋分解边界与临界溶解温度}
\noindent\textbf{题目}\par
二元液体的无量纲物质的量分数x在(0,1)，混合Gibbs能g\_mix=RT[xlnx+(1−x)ln(1−x)]+Ωx(1−x)，相互作用参数Ω=6000J/mol，温度T=300.0K且Ω视为常量。求二阶导数和两条自旋线组成、临界温度、x=0.5时均相是否局部稳定；从数学上区分spinodal与binodal，并证明自旋线何时不存在真实根。\par
\medskip\noindent\textbf{分步解答}
\begin{enumerate}
\item 计算一阶导数g′=RTln[x/(1−x)]+Ω(1−2x)。
\item 二阶导数g″=RT/[x(1−x)]−2Ω，局部稳定要求g″>0。
\item 自旋线令x(1−x)=RT/(2Ω)，即x±=[1±√(1−2RT/Ω)]/2。
\item 带入300K得xlow=0.294723517、xhigh=0.705276483。
\item 临界判别式为零时Tcrit=Ω/(2R)=360.817065K；高于Tc无内部自旋根。
\item x=0.5时g″=4RT−2Ω=-2022.64486J/mol<0，故局部不稳定；两相共存线需由两相化学势相等/公共切线决定，而不等于自旋线。
\item 验算：以独立方程将结果重新带入热力学、物料及边界条件，归一化残差为 0.000e+0、0.000e+0、0.000e+0。
\end{enumerate}
\noindent\textbf{答案摘要}\par
x\_low = 0.294723517 fraction；x\_high = 0.705276483 fraction；Tc = 360.817065 K；g\_second = -2022.64486 J/mol\par
\noindent\textbf{模型前提：}
理想正则溶液模型、无波动临界修正，忽略体积效应与Ω温度变化。\par
\medskip\hrule\medskip
\Needspace{7\baselineskip}
\subsubsection*{UC-ADV5-P20 · 反应热容差的van’t Hoff积分修正与平衡转化反算}
\noindent\textbf{题目}\par
某假想异构化A⇌B在300K的标准反应焓ΔH°=−40.00kJ/mol、标准反应熵ΔS°=−120.0J mol\textasciicircum{}-1K\textasciicircum{}-1，反应热容差ΔCp=+80.00J mol\textasciicircum{}-1K\textasciicircum{}-1在300–350K固定。理想溶液中起始纯A且体积恒定，求300K、350K的K及350K的平衡B摩尔分数。证明为什么不能用常焓van’t Hoff指数跨越50K，并以ΔG°(T)反算。\par
\medskip\noindent\textbf{分步解答}
\begin{enumerate}
\item 温度相关反应焓ΔH°(T)=ΔH300+ΔCp(T−300)。
\item 积分热容差得到反应熵ΔS°(T)=ΔS300+ΔCp ln(T/300)。
\item 标准反应Gibbs自由能ΔG°(T)=ΔH°(T)−TΔS°(T)。
\item 热力学平衡常数K(T)=exp[−ΔG°(T)/(RT)]，所以K300=4.97105157。
\item 同样计算K350=0.560679086，理想A⇌B平衡摩尔比nB/nA=K350。
\item 起始纯A且无旁路反应，xB=K350/(1+K350)=0.359253283；若假定ΔH恒定，将丢失ΔCp在焓与熵中的双重修正。
\item 验算：以独立方程将结果重新带入热力学、物料及边界条件，归一化残差为 0.000e+0、-1.110e-16、0.000e+0。
\end{enumerate}
\noindent\textbf{答案摘要}\par
K300 = 4.97105157 dimensionless；K350 = 0.560679086 dimensionless；xB350 = 0.359253283 fraction\par
\noindent\textbf{模型前提：}
理想异构化与标准态无量纲K，热容差恒定，不考虑其他温度影响。\par
\medskip\hrule\medskip
\Needspace{7\baselineskip}
\subsubsection*{UC-ADV5-P21 · Debye–Hückel活度校正下二价电极电势的偏差定量}
\noindent\textbf{题目}\par
298.15K假想二价金属电极半反应M2++2e−⇌M(s)，金属固相活度1、E°=0.3000V，分析所得自由M2+浓度0.005000mol/L、离子强度0.02000M。题设离子活度系数采用log10γ=−0.509 z²√I/(1+√I)，z=2。求γ、标准1M参照下活度比、正确电势、浓度理想近似电势及偏差；分析为什么同浓度下二价离子受活度校正更强。\par
\medskip\noindent\textbf{分步解答}
\begin{enumerate}
\item 先取√I并按题定扩展Debye–Hückel式计算log10γ=-0.252259063。
\item 活度系数γ=10\textasciicircum{}logγ=0.559423798，M2+活度比a/c°=γc/1M=0.00279711899。
\item 半反应Nernst式E=E°+(RT/2F)ln(a/c°)，实际电位=0.22447454V。
\item 若错把浓度替代活度，Eideal=E°+(RT/2F)ln(c/c°)=0.231936281V。
\item 正确电势相对理想近似偏移ΔE=(RT/2F)lnγ=-0.00746174105V。
\item 活度偏差含z²依赖（此题z=2），价态绝对值越大相互作用校正通常越显著；实际浓度更高时需超出此近似的活度模型。
\item 验算：以独立方程将结果重新带入热力学、物料及边界条件，归一化残差为 0.000e+0、1.388e-17、-3.469e-18。
\end{enumerate}
\noindent\textbf{答案摘要}\par
gamma = 0.559423798 dimensionless；activity = 0.00279711899 relative 1M；E = 0.22447454 V；error = -0.00746174105 V\par
\noindent\textbf{模型前提：}
只采用题定Debye–Hückel近似、固态金属活度为1，忽略液接势、离子尺寸高阶项与过电位。\par
\medskip\hrule\medskip
\Needspace{7\baselineskip}
\subsubsection*{UC-ADV5-P22 · 逆流换热器NTU—效能法的双流量热量核算}
\noindent\textbf{题目}\par
稳态无相变逆流换热器的热流体热容流率Ch=400W/K、入口Th,in=360K，冷流体Cc=250W/K、入口Tc,in=290K，总传热能力UA=180W/K且不随温度改变。已知逆流效能ε=[1−exp(−NTU(1−Cr))]/[1−Cr exp(−NTU(1−Cr))]。求Cr、NTU、ε、实际热交换功率、两出口温度以及两端温差，检查能量守恒和物理温度梯度。\par
\medskip\noindent\textbf{分步解答}
\begin{enumerate}
\item Cmin=250W/K、Cmax=400W/K，所以Cr=0.625。
\item NTU=UA/Cmin=180/250=0.72。
\item 将NTU、Cr代入效能关系求ε=0.452526332。
\item 最大可能换热功率Qmax=Cmin(Th,in−Tc,in)=17500W，实际Q=εQmax=7919.21082W。
\item 热出口Th,out=Th,in−Q/Ch=340.201973K；冷出口Tc,out=Tc,in+Q/Cc=321.676843K。
\item 检查热侧放热Ch(360−Tho)与冷侧吸热Cc(Tco−290)都等于Q；两端最小温差均应为正，不能越过热力学允许的温度梯度。
\item 验算：以独立方程将结果重新带入热力学、物料及边界条件，归一化残差为 -8.039e-16、-4.594e-16、0.000e+0。
\end{enumerate}
\noindent\textbf{答案摘要}\par
NTU = 0.72 dimensionless；eps = 0.452526332 fraction；Q = 7919.21082 W；Thout = 340.201973 K；Tcout = 321.676843 K\par
\noindent\textbf{模型前提：}
恒定热容流率和总UA，理想逆流稳态，无热损失与潜热。\par
\medskip\hrule\medskip
\Needspace{7\baselineskip}
\subsubsection*{UC-ADV5-P23 · 双扩散层加线性表面反应的串联传质阻力}
\noindent\textbf{题目}\par
两种串联扩散介质与界面反应等效传质膜组成复合系统，浓度势由左向右降低，左侧体积浓度0.01000mol/m³、右侧0.002000mol/m³。三段线性传递系数k1=1.00×10\textasciicircum{}-4m/s、k2=2.00×10\textasciicircum{}-4m/s、ks=5.00×10\textasciicircum{}-5m/s。题设稳态时三段通量相等，界面反应在传质网络中以线性浓度势降表达。求共同通量、三段势降、膜中两节点浓度以及总阻力，检查各级浓度非负且总势降符合外界边界条件。\par
\medskip\noindent\textbf{分步解答}
\begin{enumerate}
\item 每段的线性通量式为J=ki Δci，故Ri=1/ki、Rtotal=ΣRi。
\item 总阻力Rtotal=35000 s/m。
\item 稳态共同通量J=(cL−cR)/Rtotal=2.28571429e-7mol/(m²s)。
\item 第一、第二扩散层的浓度势降分别为0.00228571429与0.00114285714mol/m³。
\item 线性表面反应段的势降=0.00457142857mol/m³，占总驱动力比例=0.571428571。
\item 沿方向节点浓度c1=cL−Δc1=0.00771428571、c2=cL−Δc1−Δc2=0.00657142857mol/m³；所有势降之和必须等于0.008000mol/m³。
\item 验算：以独立方程将结果重新带入热力学、物料及边界条件，归一化残差为 0.000e+0、0.000e+0、0.000e+0。
\end{enumerate}
\noindent\textbf{答案摘要}\par
J = 2.28571429e-07 mol/(m2 s)；Rtotal = 35000 s/m；drop1 = 0.00228571429 mol/m3；drop2 = 0.00114285714 mol/m3；dropSurface = 0.00457142857 mol/m3\par
\noindent\textbf{模型前提：}
题设以一个独立线性阻力代表界面反应与传质，非一般真实化学反应机理。\par
\medskip\hrule\medskip
\Needspace{7\baselineskip}
\subsubsection*{UC-ADV5-P24 · 固体润湿角控制的异相形核势垒削减因子}
\noindent\textbf{题目}\par
一假想均匀球形新相核与固体平板接触，经典异相形核近似中同质球核表面能γ=0.06500J/m²，体积驱动力Δg=1.000×10\textasciicircum{}5J/m³，异相润湿接触角θ=70.00°，温度300K。球核临界半径r*=2γ/Δg、均相形核势垒Ghom*=16πγ³/(3Δg²)，异相削减因子f(θ)=(2+cosθ)(1−cosθ)²/4。求临界半径、均相与异相势垒、削减因子和异相/均相指数形核概率比，并说明θ→0、180°的极限。\par
\medskip\noindent\textbf{分步解答}
\begin{enumerate}
\item 球形经典形核自由能ΔG(r)=4πγr²−(4π/3)Δg r³，极值给r*=2γ/Δg=0.0000013m。
\item 无表面时均相势垒Ghom*=16πγ³/(3Δg²)=4.60138604e-13J。
\item 润湿角θ=70°，异相几何削减因子fθ=0.253487082。
\item 异相势垒Ghetero*=fθ Ghom=1.16639192e-13J。
\item 同一动力学前因子假设下相对指数概率比例Jhetero/Jhom≈exp[(Ghom−Ghetero)/(kBT)]，自然对数为82931870.5。
\item θ→0时fθ→0，几何上完全润湿势垒趋0；θ→180°时fθ→1，异相趋于均相，仍须注意真实表面异质和动力学前指数。
\item 验算：以独立方程将结果重新带入热力学、物料及边界条件，归一化残差为 0.000e+0、0.000e+0、0.000e+0。
\end{enumerate}
\noindent\textbf{答案摘要}\par
rstar = 1.3e-06 m；Ghom = 4.60138604e-13 J；f = 0.253487082 fraction；Ghet = 1.16639192e-13 J\par
\noindent\textbf{模型前提：}
经典球冠形核几何模型，接触角固定，忽略接触线张力与动力学前因子差异。\par
\medskip\hrule\medskip
\Needspace{7\baselineskip}
\subsubsection*{UC-ADV5-P25 · 有限两池理想溶质混合的Gibbs能下降与最大可逆功}
\noindent\textbf{题目}\par
两个不同浓度理想单一中性溶质水溶液在298.15K、固定体积下分别为V1=1.000L、c1=0.1000M，V2=2.000L、c2=0.02000M。假设溶剂与各池体积不变、溶质稀溶液的混合Gibbs能对浓度的相关部分为G=RTΣ n\_i(ln(c\_i/c°)−1)，标准c°=1M。求直接混合的最终浓度、总溶质mol、Gibbs能变化和理论可提取最大可逆非膨胀功，证明ΔG≤0并说明为什么仅用总溶质摩尔数无法算出混合功。\par
\medskip\noindent\textbf{分步解答}
\begin{enumerate}
\item 总溶质mol数n=c1V1+c2V2=0.14，总容量Vtotal=3.000L。
\item 混合后均匀浓度cf=n/(V1+V2)=0.04666666667M。
\item 理想稀溶质的标准态贡献与−RTn项在前后相同总n下严格相消。
\item 因此ΔG=RT\{n ln(cf/c°)−n1 ln(c1/c°)−n2 ln(c2/c°)\}=-104.9146445J。
\item 最大可逆非膨胀功Wrev,max=−ΔG=104.9146445J，正值代表浓差资源可利用。
\item 由xlnx的凸性或Gibbs化学势均衡可证明ΔG≤0；即使总mol相同，不同初始浓度与分池体积会给出不同混合功。
\item 最后以独立计算核对热力学力、反应动力学与守恒，归一化残差为 4.956e-17、0.000e+0、1.137e-13。
\end{enumerate}
\noindent\textbf{答案摘要}\par
cf = 0.0466666667 mol/L；n = 0.14 mol；deltaG = -104.914644 J；maxWork = 104.914644 J\par
\noindent\textbf{模型前提：}
理想稀溶液且体积相加，不计算实际溶剂化与膜不可逆耗散，理想标准态及温度固定。\par
\medskip\hrule\medskip
\Needspace{7\baselineskip}
\subsubsection*{UC-ADV5-P26 · 球形多孔催化颗粒一级反应的效能因子与表面浓度}
\noindent\textbf{题目}\par
等温球形多孔催化剂颗粒半径Rp=2.000mm、有效扩散系数De=1.00×10\textasciicircum{}-9m²/s，颗粒内反应物按一级不可逆反应−r=k c、k=0.004000s\textasciicircum{}-1，表面反应物浓度cs=0.01000mol/m³由外部充分搅拌维持。球对称稳态满足De(1/r²)d/dr[r²dc/dr]−kc=0，中心通量0。求Thiele模数、催化效能因子η=3(φ cothφ−1)/φ²、颗粒中心/表面浓度比及每颗粒实际反应量与无内扩散限制反应量之比；证明η<1并给出φ→0极限。\par
\medskip\noindent\textbf{分步解答}
\begin{enumerate}
\item 对球对称反应扩散方程，令λ=√(k/De)，无量纲Thiele模数φ=Rpλ=4。
\item 保持中心有限的解析解为c(r)/cs=[Rp sinh(λr)]/[r sinhφ]，r→0需取极限。
\item 中心/表面浓度比c(0)/cs=φ/sinhφ=0.1465742813。
\item 将球内k c(r)对体积积分并与无扩散限制kcs的体积总速率相比，得η=3(φ cothφ−1)/φ²。
\item 对参数计算η=0.5630033628，由于内部浓度均小于或等于cs，0<η<1。
\item φ→0时cothφ=1/φ+φ/3+O(φ³)，因此η→1；φ大时η≈3/φ，解释强扩散限制下体积利用效率下降。
\item 最后以独立计算核对热力学力、反应动力学与守恒，归一化残差为 0.000e+0、0.000e+0、0.000e+0。
\end{enumerate}
\noindent\textbf{答案摘要}\par
phi = 4 dimensionless；center\_ratio = 0.146574281 fraction；eta = 0.563003363 fraction；boundary\_factor = 3.0026846 dimensionless\par
\noindent\textbf{模型前提：}
球形均匀多孔颗粒、有效D恒定、一级不可逆等温反应，外部膜阻力已忽略。\par
\medskip\hrule\medskip
\Needspace{7\baselineskip}
\subsubsection*{UC-ADV5-P27 · 全混流CSTR的一阶反应与瞬态响应时间常数}
\noindent\textbf{题目}\par
一台恒温完全混合连续搅拌釜CSTR中进行假想单一液相一级反应A→B，进料浓度CA,in=0.1000M，停留时间τ=2.000min，一级动力学常数k=0.4000min\textasciicircum{}-1，体积与流量恒定。起始釜内没有A，t=0起接通恒定进料。推导瞬态A浓度微分方程并求任意时刻C\_A(t)，计算最终稳态浓度、稳态转化率、达到稳态浓度90\%的时间与当时产物生成速率，证明混合釜稳态与同停留时间PFR转化率不同。\par
\medskip\noindent\textbf{分步解答}
\begin{enumerate}
\item 完全混合釜物料守恒V dC/dt=F(Cin−C)−Vk C，化为dC/dt=Cin/τ−(1/τ+k)C。
\item 定义衰减特征率λ=1/τ+k=0.9 min\textasciicircum{}-1，积分因子解C(t)=Css[1−e\textasciicircum{}(−λt)]。
\item 稳态浓度Css=Cin/(1+kτ)=0.05555555556M。
\item 稳态转化率X=kτ/(1+kτ)=0.4444444444；PFR一级模型为1−e\textasciicircum{}(−kτ)，不会与CSTR相同。
\item 达到90\%稳态浓度需要t90=ln10/λ=2.558427881min。
\item 此时C\_A=0.9Css，产物体积生成速率rB=kC\_A=0.02M/min；A的瞬态积累须按完整CSTR守恒而非直接套用批反应。
\item 最后以独立计算核对热力学力、反应动力学与守恒，归一化残差为 0.000e+0、0.000e+0、0.000e+0。
\end{enumerate}
\noindent\textbf{答案摘要}\par
CAss = 0.0555555556 mol/L；conversion = 0.444444444 fraction；lambda = 0.9 min\textasciicircum{}-1；t90 = 2.55842788 min\par
\noindent\textbf{模型前提：}
完全混合单一液相、恒流量恒体积无旁路与死区，一级不可逆。\par
\medskip\hrule\medskip
\Needspace{7\baselineskip}
\subsubsection*{UC-ADV5-P28 · Onsager互易耦合热质传递的熵产生正定性检验}
\noindent\textbf{题目}\par
在线性不可逆热力学近似下，两类标量热力学通量J1、J2和共轭广义热力学力X1、X2满足Onsager互易关系[J1,J2]\textasciicircum{}T=L[X1,X2]\textasciicircum{}T，L=[[3.000,0.8000],[0.8000,2.000]]（按给定归一化单位），两个外驱动力X1=0.2000、X2=−0.1000。求两通量、总无量纲熵产生率σ=J1X1+J2X2、交叉项对熵产生的贡献；证明L为正定矩阵并解释两通量为何可能在某条件下逆自身直接力方向。\par
\medskip\noindent\textbf{分步解答}
\begin{enumerate}
\item Onsager矩阵按题设L12=L21=0.8000，满足互易关系。
\item 通量J1=3X1+0.8X2=0.52，J2=0.8X1+2X2=-0.04。
\item 总熵产生σ=X\textasciicircum{}T L X=J1X1+J2X2=0.108，应非负。
\item 对角贡献L11X1²+L22X2²=0.14，交叉贡献2L12X1X2=-0.032。
\item 矩阵正定充分必要条件L11>0且det L=L11L22−L12²=5.36>0，因此对任何非零X都有σ>0。
\item 虽然总熵产生必为正，单一通量可由交叉项主导而与该通量自己的驱动力方向相反，并通过耦合另一通量而不违反第二定律。
\item 最后以独立计算核对热力学力、反应动力学与守恒，归一化残差为 0.000e+0、0.000e+0、5.204e-18。
\end{enumerate}
\noindent\textbf{答案摘要}\par
J1 = 0.52 arb；J2 = -0.04 arb；sigma = 0.108 arb；det = 5.36 dimensionless\par
\noindent\textbf{模型前提：}
线性近均衡、两个力通量定义确保σ=ΣJiXi，L对称正定且各单位已标准化。\par
\medskip\hrule\medskip
\Needspace{7\baselineskip}
\subsubsection*{UC-ADV5-P29 · Langmuir–Hinshelwood双底物竞争覆盖的反应速率极值}
\noindent\textbf{题目}\par
假想多相表面催化反应A(ads)+B(ads)→P遵循准平衡Langmuir竞争占位，两个底物分别的吸附常数KA=100L/mol、KB=50.00L/mol，溶液自由浓度[A]=0.02000M、[B]=0.03000M，表面占位分数θA=KA[A]/(1+KA[A]+KB[B])，θB同理，表面反应速率r=kθAθB，k=0.01000mol/(g s)。求两种覆盖率、空位分数、反应速率；当B固定时严格求使速率最大的A浓度及最大速率，论证过量A如何反而抑制反应。\par
\medskip\noindent\textbf{分步解答}
\begin{enumerate}
\item 设α=KA[A]=2、β=KB[B]=1.5，竞争占位分母D=1+α+β=4.5。
\item 两种占位率θA=α/D=0.4444444444、θB=β/D=0.3333333333，空位率θ*=1/D=0.2222222222。
\item 速率r=kθAθB=0.001481481481 mol/(g s)。
\item 固定B使β恒定，速率正比于αβ/(1+β+α)²。
\item 对α求导为正比于(1+β−α)/(1+β+α)³，唯一极大值αopt=1+β=2.5。
\item 因此[A]opt=(1+KB[B])/KA=0.025M，最大速率rmax=kβ/[4(1+β)]=0.0015；A过多会占满位点并挤走B。
\item 最后以独立计算核对热力学力、反应动力学与守恒，归一化残差为 -1.110e-16、0.000e+0、0.000e+0。
\end{enumerate}
\noindent\textbf{答案摘要}\par
thetaA = 0.444444444 fraction；thetaB = 0.333333333 fraction；rate = 0.00148148148 mol/(g s)；Aopt = 0.025 mol/L\par
\noindent\textbf{模型前提：}
表面竞争Langmuir单层占位，表面反应为双位点碰撞且不改变准平衡吸附，不考虑位点邻接相关。\par
\medskip\hrule\medskip
\Needspace{7\baselineskip}
\subsubsection*{UC-ADV5-P30 · 化学电池有限可逆反应的反应商、自由能与电动势}
\noindent\textbf{题目}\par
298.15K假想两电子电池总反应R⇌P，标准态下平衡常数K=100.0，当前理想活动度反应商Q=0.1500，n=2，标准态理想反应Gibbs能ΔG°=−RTlnK。某小段反应正向推进0.01000mol时，把电位在这一小段视为瞬时初值近似，求当前ΔG反应摩尔自由能、电动势E、按当前电势估算的小段电功及标准电势，并明确为何较大有限转化必须积分E(ξ)dQ而不能一直用初值。\par
\medskip\noindent\textbf{分步解答}
\begin{enumerate}
\item 反应标准Gibbs能ΔG°=−RT ln K=-11416.019J/mol。
\item 在当前反应商Q=0.15时ΔG=ΔG°+RTlnQ=RTln(Q/K)=-16118.89793J/mol。
\item 对应瞬时可逆电动势E=−ΔG/(nF)=0.08353030234V。
\item 标准电动势E°=RTlnK/(nF)=0.05915934969V，二者相差(−RT/nF)lnQ。
\item 小段0.01000mol反应转移电量dQ≈nF×0.01000=1929.706642C，初值电势估算的可逆电功≈E dQ=161.1889793J。
\item 当反应商随转化显著改变，E不是常数，应对电势随进度积分∫E(ξ)nF dξ，不能以当前初值乘整个容量。
\item 最后以独立计算核对热力学力、反应动力学与守恒，归一化残差为 0.000e+0、0.000e+0、0.000e+0。
\end{enumerate}
\noindent\textbf{答案摘要}\par
deltaG = -16118.8979 J/mol；E = 0.0835303023 V；Estandard = 0.0591593497 V；approx\_work = 161.188979 J\par
\noindent\textbf{模型前提：}
电池为题定理想化两电子体系，当前电势用于小段线性近似而非完整放电能量；活动度反应商无量纲。\par
\medskip\hrule\medskip
\end{document}
